% concatenated from main.tex
%lono%%%%%%%%%%%%%%%%%%%%%%%%%%%%%%%%%%%%%%%%%%%%%%%%%%%%%%%%%%%%%%%%%%%%%%%%%%%%%%%
%2345678901234567890123456789012345678901234567890123456789012345678901234567890
%        1         2         3         4         5         6         7         8

\documentclass[letterpaper, 10 pt, conference]{ieeeconf}  % Comment this line out
                                                          % if you need a4paper
%\documentclass[a4paper, 10pt, conference]{ieeeconf}      % Use this line for a4
                                                          % paper

\IEEEoverridecommandlockouts                              % This command is only
                                                          % needed if you want to
                                                          % use the \thanks command
\overrideIEEEmargins
% See the \addtolength command later in the file to balance the column lengths
% on the last page of the document
\usepackage{cite}
\usepackage{amsmath,amssymb,amsfonts}
\usepackage{algorithmic}
\usepackage{graphicx}
\usepackage{textcomp}
\usepackage{xcolor}
\usepackage[scr]{rsfso}
\usepackage{authblk}
\usepackage{float}
\usepackage{soul}
\usepackage{hyperref}
\usepackage{url}
\usepackage{float}
\usepackage{titlesec}
\usepackage{booktabs} % For professional table formatting
\usepackage{array}    % For column alignment
\usepackage{multirow} % For multirow cells
\usepackage{siunitx}  % For aligning numbers
\usepackage{caption}  % For caption control
\usepackage{makecell} % For line breaks in cells
\usepackage{mathtools}
\usepackage{etoolbox} 
\usepackage{graphicx}
\usepackage{svg}
\usepackage{subcaption}
\usepackage{bm}


\usepackage{hyperref}
\hypersetup{
    colorlinks=true,
    linkcolor=blue,
    citecolor=red,
    urlcolor=magenta
}
\usepackage{lmodern}
\usepackage{titlesec}

\newcommand{\norm}[1]{\left\lVert#1\right\rVert}

\newtheorem{theorem}{Theorem}
\newtheorem{definition}[theorem]{Definition}
\newtheorem{assumption}[theorem]{Assumption}
\newtheorem{lemma}[theorem]{Lemma}
\newtheorem{corollary}[theorem]{Corollary}
\newtheorem{example}[theorem]{Example}

\newcommand{\qedsymbol}{$\blacksquare$}

\usepackage{xcolor}
\newcommand\todo[1]{\textcolor{red}{#1}}

% Define the counter per section
\newcounter{parcount}[section]
\renewcommand{\theparcount}{\textit{(\alph{parcount})}}

% Custom paragraph-like label
\newcommand{\myparagraph}[1]{%
  \par\refstepcounter{parcount}%
  \noindent\hspace*{1em}\theparcount \textit{ #1}:
}

% Ensure counter resets with section
\makeatletter
\@addtoreset{parcount}{section}
\makeatother

\newcommand{\rplus}[0]{\mathbb{R}_+}

\begin{document}
\title{
Delay compensation of multi-input distinct delay \\ nonlinear systems via
neural operators}

\author[1]{Filip Bajraktari}
\author[2]{Luke Bhan\thanks{Correspondence to: lbhan@ucsd.edu}}
\author[2]{Miroslav Krstic}
\author[2]{Yuanyuan Shi}

\affil[1]{University of Belgrade, Serbia}
\affil[2]{University of California San Diego, La Jolla, CA, USA}

% Reduce spacing between author lines (IEEE-style)
\makeatletter
\patchcmd{\@maketitle}{\AB@authlist}{\AB@authlist\vspace{-1.5em}}{}{}
\makeatother

% Title
\maketitle
\begin{abstract}
In this work, we present the first stability results for approximate predictors in multi-input nonlinear systems with distinct actuation delays. We show that if the predictor approximation satisfies a uniform (in time) error bound, semi-global practical stability is correspondingly achieved.  For such approximators, the required uniform error bound depends on the desired region of attraction and the number of control inputs in the system. The result is achieved through transforming the delay into a transport PDE and conducting analysis on the coupled ODE-PDE cascade. To highlight the viability of such error bounds, we demonstrate our results on a class of approximators - neural operators - showcasing sufficiency for satisfying such a universal bound both theoretically and in simulation on a mobile robot experiment. 
\vspace*{\baselineskip}

\noindent \textbf{Keywords:} Nonlinear Systems, Delay Systems, Deep Learning, Neural Operators
\end{abstract}

\section{Introduction} \label{sec:introduction}
We study the multi-input, distinct delay, nonlinear system governed by
\begin{equation}
    \dot{X}(t) = f \bigg( X(t), U_1(t-D_1), \dots, U_m(t-D_m) \bigg), \label{eq:system}
\end{equation}
where $X \in \mathbb{R}^n$ is the system state, $U_1, \dots, U_m$ are scalar control inputs, $f: \mathbb{R}^n \times \mathbb{R}^m \rightarrow \mathbb{R}^n$ is a locally Lipschitz vector field satisfying $f(0, \dots,0) = 0$, and $D_1, \dots, D_m$ are (potentially distinct) input delays satisfying (without the loss of generality) $0 < D_1 \leq \dots \leq D_m$. This system commonly appears across a range of applications, including telerobotics \cite{9727201}, connected and automated vehicles \cite{samii2025experimentalimplementationvalidationpredictorbased}, networking systems \cite{doi:10.1137/140980570}, and many more \cite{Bekiaris-Liberis_Krstic:2016}.

To compensate for the delays, the most popular approach for nonlinear systems, termed ``predictor feedback", was developed in \cite{Krstic:2010}. The key idea of predictor feedback designs is to apply known delay-free, stabilizing controllers with a future prediction of the state; thus, when the control law arrives at the system, it aligns with the system's true state. Consequently, for any predictor feedback design, we  require the existence of a set of globally asymptotically stabilizing delay-free control laws:
\begin{assumption}\label{assumption:delay-free-control-laws}
There exists a set of controllers $U_i(t) = \kappa_i \big( X(t)\big)$ with $\kappa_i \in C^1(\mathbb{R}^n;\mathbb{R})$, $i \in [m]$, such that the delay-free system $\dot{X}(t) = f \bigg( X(t), U_1(t), \dots, U_m(t) \bigg)$ is globally asymptotically stable.
\end{assumption}
\noindent In multi-input multi-delay systems, a key challenge naturally arises: with multiple inputs and associated delays, it becomes difficult to coordinate the predictors so that all $m$ control inputs arrive with the correct corresponding state predictions. This was resolved for \emph{exact} predictor feedback designs in \cite{Bekiaris-Liberis_Krstic:2016} with predictors defined as
\allowdisplaybreaks
\begin{subequations}\label{eqs:P_system}
    \begin{align}
        P_1(t) =& X(t) + \int_{t-D_1}^t f \bigg(P_1(\theta), U_1(\theta), U_2(\theta - D_{2,1}),\nonumber \\ 
        &   \dots, U_m(\theta-D_{m,1}) \bigg) d\theta \label{eq:P1(t)} \\
        P_2(t) =& P_1(t) + \int_{t-D_{2,1}}^t f \bigg( P_2(\theta), \kappa_1(P_2(\theta)), U_2(\theta), \nonumber \\& U_3(\theta - D_{3,2}),\dots, U_m(\theta-D_{m,2}) \bigg) d\theta \label{eq:P2(t)} \\
        &\hspace{3cm} \vdots \nonumber \\
        P_m(t) =& P_{m-1}(t) + \int_{t-D_{m,m-1}}^t f \bigg( P_m(\theta), \kappa_1(P_m(\theta)), \nonumber \\ & \kappa_2(P_m(\theta)), \dots, \kappa_{m-1}(P_m(\theta)), U_m(\theta) \bigg) d\theta, \label{eq:Pm(t)}
    \end{align}
\end{subequations}
\allowdisplaybreaks
with initial conditions given by 
\begin{subequations}\label{eqs:P_init_system}
\begin{alignat}{2}
  P_1(\theta) 
   = &X(0) + \int_{-D_1}^\theta f\big( P_1(s), U_1(s), U_s(&&\!\!\!\!s- D_{2,1}), \nonumber  \\
  &\dots, U_m(s - D_{m,1}) \big) ds&&  \\ 
  & && \hspace{-0.5cm} -D_1 \leq \theta \leq 0,\nonumber  \\
P_2(\theta) 
    =& P_1(0) + \int_{-D_{2,1}}^\theta f\big( P_2(s), \nonumber \kappa_1(P_2(s)&&\!\!\!\!), U_2(s), 
 \\&U_3(s - D_{3,2}), \dots, U_m(s - D_{m,2}) &&\!\!\!\!\big) ds \\ 
  & && \hspace{-0.7cm} -D_{2,1} \leq \theta \leq 0, \nonumber \\
  & \hspace{3.5cm} \vdots \nonumber  && \\ 
P_m(\theta) =& P_{m-1}(0) + \int_{-D_{m,m-1}}^\theta f \bigg( P_m(s), \kappa_1&&(P_m(s)), \nonumber \\ &\kappa_2(P_m(s)), \dots, \kappa_{m-1}(P_m(s)), U_m&&(s) \bigg) ds \label{eq:Pm(theta)}\,, \\
            & && \hspace{-1.3cm}  -D_{m,m-1} \leq \theta \leq 0\,, \nonumber
    \end{alignat}
\end{subequations}
where $D_{i, j}$ represents the difference between delays, $D_i - D_j$. 

While the predictors in \eqref{eqs:P_system}, \eqref{eqs:P_init_system} address the issue of distinct delays, the original global stability result \cite[Theorem 1]{Bekiaris-Liberis_Krstic:2016} only applies to exact implementations—which are never feasible without closed-form solutions. Thus, as implicit ODEs, implementing the predictors numerically requires both fixed-point iteration and integral discretization \cite{Predictor_Feedback_For_Delay_Systems}. This poses two main challenges: (i) discretization errors grow with step size and delay, potentially destabilizing the feedback, and (ii) convergence of the fixed-point iteration is not guaranteed, especially for stiff systems. To address this, we make two contributions: (1) a stability result for approximate predictors with uniform error bounds, yielding a semi-global region of attraction; and (2) a neural network-based predictor approximation that significantly accelerates computation compared to traditional methods. For the context, we begin with an overview of recent works in predictor feedback and learning-based approximator designs.

\myparagraph{Delay compensation in nonlinear systems}
Since the Smith predictor \cite{smith1957closer}, predictor feedback has been a foundational tool for delay compensation in nonlinear systems. It has been successfully applied in diverse settings, including delay-adaptive control \cite{6704718}, PDE-based control \cite{9528941}, and event-triggered control \cite{NOZARI2020108754}. In this work, we focus on the predictor feedback designs introduced in \cite{Krstic:2010, Nonlinear_Control_Under_Nonconstant_Delay}, which model input delays as transport PDEs and apply backstepping transformations to guarantee stability. This PDE-based approach has enabled advances in controlling systems with delays across several domains, including connected and automated vehicles (CAVs) \cite{samii2025experimentalimplementationvalidationpredictorbased}, biological systems \cite{6704718}, and additive manufacturing \cite{7921611}. For a comprehensive overview of predictor feedback methods, see \cite{Deng10092022}.

\myparagraph{Learning-based approximations in control}
For decades, neural networks hve been approximating various computationally intractable portions of control. Perhaps most famously, these types of approaches have been long studied for the Hamilton-Jacobi-Bellman equations \cite{4359214}, as well as ODE approximations \cite{SONTAG1998151}, differential games \cite{pmlr-v283-sharpless25a}, and even in the context of operator approximations in PDEs \cite{Luke_gain:2024, KRSTIC2024111649}. They were recently introduced for predictors in \cite{Luke_Peijia:2025} and have since been extended for unknown delays in \cite{CDC}. We emphasize that there are many many more results not present here, but due to space constraints, must be omitted.

\myparagraph{Notation}
We use $\mathbb{R}_+$ to denote the set of positive reals. For an $n$-vector, the norm $|\cdot|$ denotes the usual Euclidean norm. For a function $u: [0,1]\times\mathbb{R} \rightarrow \mathbb{R}$, we denote by $\|u(t)\|_{L^\infty}$ the supremum spatial norm ($\|u(t)\|_{L^\infty} = \sup_{x \in [0,1]}|u(x,t)|$). We use the PDE notation $u_x(x,t)=\frac{\partial u}{\partial x}(x,t)$ to denote derivatives. We use $C^n([t-D,t];\mathbb{R}^m)$ to denote the set of functions with continuous $n$ derivatives. For brevity, when applicable we use $i \in [m]$ to represent the iteration $i \in \{1, 2, 3, \cdots , m\}$.

\section{An operator theoretic view of predictors} \label{sec:predictor-operator}
In this section, we begin by establishing a formal mapping of the solution to the implicit predictors in \eqref{eqs:P_system}, \eqref{eqs:P_init_system}. This perspective, first introduced in \cite{Luke_Peijia:2025}, enables us to view the solution to the implicit predictor ODE as a mapping of the current state and $m$ different control histories ($U_i$, $i \in [m]$) into an estimate of the system at $X(t+D_i)$. However, this definition was originally introduced for a single-input scalar system and thus requires extension to the multi-input multi-delay case. Henceforth, we begin by introducing the definition of $m$ different predictor operators
\begin{definition}\label{def:delay_dependent_predictor_operators}
(\textbf{Predictor operators for multi-input multi-delay systems}) Let $Q \in \mathbb{R}^n$, and for each $i \in [m]$, let $U_i \in C^1([0,1];\mathbb{R})$. Let $\varphi \in \mathbb{R}_+$. Then, define the set of predictor operators $\{\mathcal{P}_1, ..., \mathcal{P}_m\}$ mapping from $\mathbb{R}^n \times C^1([0,1];\mathbb{R})^{m-i+1} \times \mathbb{R}_+$ to $C^1([0,1];\mathbb{R}^n)$ taking
 inputs $(Q,U_i,\dots,U_m, \varphi)$ and returning a function $P_i \in C^1([0,1];\mathbb{R}^n)$ where the output $P_i$ satisfies the implicit differential equation
\begin{align}
    0 = P_i(s) - Q& - \varphi \int_0^s f(P_i(\theta), \kappa_1(P_i(\theta)), \dots, \kappa_{i-1}(P_i(\theta)),\nonumber \\ 
    &U_i(\theta), \dots, U_m(\theta)) d\theta , \quad s \in [0,1], \label{eq:predictor_def_Pi}
\end{align}
where $\kappa_i$, $i \in [m-1]$, are a set of $C^1$ functions from $\mathbb{R}^n \to \mathbb{R}$ that are assumed to be apriori known as static parameters as in Assumption \ref{assumption:delay-free-control-laws}.
\end{definition}
Notice that, by definition, the predictor operator yields the solution to \eqref{eqs:P_system} in the following sense
\begin{subequations}
\begin{align}
    P_1(t) =& \mathcal{P}_1(X(t), T_{D_1,0}(t)U_1, ..., \nonumber \\ & \quad  T_{D_m,D_{m-1}}(t)U_m, D_1)(1), \\
    P_2(t) =& \mathcal{P}_2(P_1(t), T_{D_{2,1},0}(t)U_2, ..., \nonumber \\ & \quad T_{D_{m,1},D_{m,2}}(t)U_m, D_{2,1})(1), \\
    & \hspace{2.5cm} \vdots \nonumber \\
    P_m(t) =& \mathcal{P}_m(P_{m-1}(t), T_{D_{m,m-1},0}(t)U_m, D_{m,m-1})(1),
\end{align}
\end{subequations}
where $T_{\varphi_i,\varphi_j}$, $i, j \in [m]$, $i \geq j$, are the operators that yield the control history of $U$ defined as
\begin{align}
    (T_{\varphi_i,\varphi_j}(t)U) := U(t-\varphi_i+(\varphi_i-\varphi_j)x), \quad x \in [0,1).
\end{align}



The choice to use $m$ separate predictor operators instead of one large operator offers is both deliberate and advantageous. It allows training $m$ separate neural networks on specific datasets rather than a single monolithic model, simplifying training and enabling parallel inference across compute nodes for significant speedups. The downside, however, is that this modular design, causes output discontinuities at operator boundaries as 
$P_{j}(\cdot)(1) \neq P_{j+1}(\cdot)(0)$, 
$j \in [m-1]$. 
Consequently, this is only a theoretical inhibition as practicality, one can enforce continuity in implementation, setting $P_{j+1}(\cdot)(0) = P_{j}(\cdot)(1)$ or following the nested operator design \cite{D2EE04204E}. Moreover, by design, only the minimal information for each predictor—controlled by the deliberate shift $T_{\varphi_1, \varphi_2}(t)$—is needed, reducing computational load.  Lastly, scaling predictors on $s \in [0, 1]$ is intentional, allowing reuse without retraining when the input context $\varphi$ changes and thus, enables scaling across different delays.

\subsection{Continuity of the predictor operators} \label{subsec:continuity}
 In anticipation of approximating the predictor operators in Section \ref{subsec:neural-op-approximators}, we showcase that the operator is Lipschitz continuous.
 To do so, it is standard in the predictor feedback literature \cite{Predictor_Feedback_For_Delay_Systems, Luke_Peijia:2025} to assume Lipschitz continuity of the dynamics over an arbitrary large, but compact region of the global state space. 

 
\begin{assumption}\label{assumption:lipschitz}
Let $f(X,U_1,\dots,U_m)$ be as in \eqref{eq:system} and $\mathcal{X} \subset \mathbb{R}^n$, $\mathcal{U} \subset \mathbb{R}$ be two compact sets. Then, assume there exists a constant $C_f > 0$ such that $f$ satisfies the Lipschitz condition
\begin{align}
    \big| f(x_1,u_1,\dots,u_m&)-f(x_2,u_{m+1},\dots,u_{2m}) \big|\nonumber  &&\\ \hspace{1cm} \leq& C_f \bigg( |x_1-x_2| + \sum_{i=1}^m |u_i - u_{m+i}| \bigg),
\end{align}
for all $x_1,x_2 \in \mathcal{X}$ and $u_1,\dots,u_{2m} \in \mathcal{U}$.
\end{assumption}

Now, we prove the Lipschitz continuity of the predictor operator in the following lemma.

\begin{lemma}\label{lemma:predictor-continuity}
Let $\mathcal{X}$ and $\mathcal{U}$ be compact sets and $C_f$ be the Lipschitz constant as in Assumption~\ref{assumption:lipschitz}. Let $X_1, X_2 \in \mathcal{X}$, control functions $U_j, U_{m+j} \in C^1([0,1];\mathcal{U})$ with $j \in [m]$, and positive real numbers $\varphi_1, \varphi_2 \in \Phi \subset \mathbb{R}_+$ where $\Phi$ is compact. Further, let $\overline{\mathcal{X}}, \overline{\mathcal{U}}$, and $\overline{\Phi}$ be upper bounds for each of their underlying sets. Then, each predictor operator $\mathcal{P}_i$, $i \in [m]$, from Definition~\ref{def:delay_dependent_predictor_operators} satisfies the following Lipschitz condition
\begin{align}
    \bigg\| \mathcal{P}_i&(X_1,U_i,\dots,U_m,\varphi_1) - \mathcal{P}_i(X_2,U_{m+i},\dots,U_{2m},\varphi_2) \bigg\|_{L^\infty} \nonumber  \\
    \leq& C_{\mathcal{P}} \Biggl( |X_1 - X_2| + \sum_{j=i}^m \bigg\| U_j - U_{m+j} \bigg\|_{L^\infty} 
    + |\varphi_1 - \varphi_2| \Biggr), \label{eq:predictor_lipschitz_inequality}
\end{align}
with Lipschitz constant
\begin{align}
    C_{\mathcal{P}} :=& \max \bigg( 1, \Xi, \overline{\Phi} C_f \bigg) e^{\overline{\Phi} C_{\kappa}}, \\ 
    \Xi  :=& m C_f \overline{\mathcal{U}} + C_{\kappa} \bigg( \overline{\mathcal{X}} + m \overline{\Phi} C_f \overline{\mathcal{U}} \bigg) e^{\overline{\Phi} C_{\kappa}}, \\
    C_{\kappa}  :=& C_f \bigg( 1 + \sum_{i=1}^m C_{\kappa_i} \bigg), \label{eq:c_kappa}
\end{align}
where $C_{\kappa_i}$ is Lipschitz constant of each control law $\kappa_i$.
\end{lemma}
\begin{proof}
Define the shorthand notation
\begin{align}
    P_i^{(1)}(s) &:= \mathcal{P}_i(X_1,U_i,\dots,U_m, \varphi_1)(s) \\
    P_i^{(2)}(s) &:= \mathcal{P}_i(X_2,U_{m+i},\dots,U_{2m}, \varphi_2)(s),
\end{align}
for all $s \in [0,1]$. Before we proceed to prove the main part of the lemma, we show that each predictor is uniformly bounded in its domain. First, by definition, Assumption's \ref{assumption:delay-free-control-laws}, \ref{assumption:lipschitz} and the compactness of $\mathcal{X}$, we have
\begin{align}
    \bigg| P_i^{(1)}(s) \bigg| \leq& |X_1| + \varphi_1 \int_{0}^s \bigg| f \bigg( P_i^{(1)}(\theta), \kappa_1 \bigg( P_i^{(1)}(\theta) \bigg)\nonumber \\ 
    &\quad , \dots, \kappa_{i-1} \bigg( P_i^{(1)}(\theta) \bigg), U_i(\theta), \dots, U_m(\theta) \bigg) \bigg| d\theta \nonumber \\
        % \leq& |X_1| + \varphi_1 \int_{0}^s C_f \Biggl( \bigg| P_i^{(1)}(\theta) \bigg| + \sum_{j=1}^{i-1} \bigg| \kappa_j \bigg( P_i^{(1)}(\theta) \bigg) \bigg| \nonumber \\ & + \sum_{j=i}^m \bigg| U_j(\theta) \bigg| \Biggr) d\theta \nonumber \\
        \leq& |X_1| + \varphi_1 \int_{0}^s C_f \Biggl( \bigg| P_i^{(1)}(\theta) \bigg| + \sum_{j=1}^{i-1} C_j \bigg| P_i^{(1)}(\theta) \bigg| \nonumber \\ & + \sum_{j=i}^m \bigg| U_j(\theta) \bigg| \Biggr) d\theta \nonumber \\
        \leq& |X_1| + \varphi_1 C_f \sum_{j=i}^m \| U_j \|_{L^\infty} \nonumber \\ & + \varphi_1 C_{\kappa} \int_{0}^s \bigg| P_i^{(1)}(\theta) \bigg| d\theta.
\end{align}
Applying Gronwall's inequality yields
\begin{align}
    \bigg| P_i^{(1)}(s) \bigg| \leq& \bigg( \overline{\mathcal{X}} + m \overline{\Phi} C_f \overline{\mathcal{U}} \bigg) e^{\overline{\Phi} C_{\kappa}}.
\end{align}
Then, we have
\begin{align}
    \bigg| P_i^{(1)}&(s) - P_i^{(2)}(s) \bigg| \nonumber \\ 
        \leq& |X_1 - X_2|  + \int_{0}^s \bigg| \varphi_1 f \bigg[ P_i^{(1)}(\theta), \kappa_1 \bigg( P_i^{(1)}(\theta) \bigg),\nonumber \\ 
            & \quad \dots, \kappa_{i-1} \bigg( P_i^{(1)}(\theta) \bigg), U_i(\theta), \dots, U_m(\theta) \bigg] \nonumber \\
            & \quad - \varphi_2 f \bigg[ P_i^{(2)}(\theta), \kappa_1 \bigg( P_i^{(2)}(\theta) \bigg), \dots, \kappa_{i-1} \bigg( P_i^{(2)}(\theta) \bigg), \nonumber \\ 
            & \quad U_{m+i}(\theta), \dots, U_{2m}(\theta) \bigg] \bigg| d\theta \nonumber \\
        \leq& |X_1 - X_2|+ |\varphi_1 - \varphi_2| \int_{0}^s \bigg| f \bigg( P_i^{(1)}(\theta), \kappa_1 \bigg( P_i^{(1)}(\theta) \bigg),\nonumber \\ 
            & \quad \dots, \kappa_{i-1} \bigg( P_i^{(1)}(\theta) \bigg), U_i(\theta), \dots, U_m(\theta) \bigg) \bigg| d\theta \nonumber \\
            &\quad + \varphi_2 \int_{0}^s \Bigl| f \bigg( P_i^{(1)}(\theta), \kappa_1 \bigg( P_i^{(1)}(\theta) \bigg), \dots,\nonumber \\ 
            & \quad \kappa_{i-1} \bigg( P_i^{(1)}(\theta) \bigg), U_i(\theta), \dots, U_m(\theta) \bigg) \nonumber \\
            &\qquad - f \bigg( P_i^{(2)}(\theta), \kappa_1 \bigg( P_i^{(2)}(\theta) \bigg), \dots, \kappa_{i-1} \bigg( P_i^{(2)}(\theta) \bigg), \nonumber \\ 
            & \quad U_{m+i}(\theta), \dots, U_{2m}(\theta) \bigg) \Bigr| d\theta \nonumber \\
        % \leq& |X_1 - X_2| \nonumber + |\varphi_1 - \varphi_2| C_f \sum_{j=i}^m \| U_j \|_{L^\infty} \nonumber \\ 
        %     & \quad + |\varphi_1 - \varphi_2| C_{\kappa} \int_{0}^s \bigg| P_i^{(1)}(\theta) \bigg| d\theta \nonumber \\
        %     &\quad + \varphi_2 \int_{0}^s C_f \Biggl( \bigg| P_i^{(1)}(\theta) - P_i^{(2)}(\theta) \bigg| \nonumber \\
        %     &\quad  + \sum_{j=1}^{i-1} \bigg| \kappa_j \bigg( P_i^{(1)}(\theta) \bigg) - \kappa_j \bigg( P_i^{(2)}(\theta) \bigg) \bigg| \nonumber \\ 
        %     & \quad + \sum_{j=i}^m \bigg| U_j(\theta) - U_{m+j}(\theta) \bigg| \Biggr) d\theta \nonumber \\
        \leq& |X_1 - X_2| \nonumber \\
            &\quad + |\varphi_1 - \varphi_2| \bigg[ m C_f \overline{\mathcal{U}} + C_{\kappa} \bigg( \overline{\mathcal{X}} + m \overline{\Phi} C_f \overline{\mathcal{U}} \bigg) e^{\overline{\Phi} C_{\kappa}} \bigg] \nonumber \\
            &\quad + \varphi_2 \int_{0}^s C_f \Biggl( \bigg| P_i^{(1)}(\theta) - P_i^{(2)}(\theta) \bigg| \nonumber \\
            &\quad  + \sum_{j=1}^{i-1} C_j \bigg( P_i^{(1)}(\theta) - P_i^{(2)}(\theta) \bigg) \nonumber \\ 
            & \quad + \sum_{j=i}^m \bigg| U_j(\theta) - U_{m+j}(\theta) \bigg| \Biggr) d\theta \nonumber \\
        \leq& |X_1 - X_2| \nonumber \\
            &\quad + |\varphi_1 - \varphi_2| \bigg[ m C_f \overline{\mathcal{U}} + C_{\kappa} \bigg( \overline{\mathcal{X}} + m \overline{\Phi} C_f \overline{\mathcal{U}} \bigg) e^{\overline{\Phi} C_{\kappa}} \bigg] \nonumber \\
            &\quad + \varphi_2 C_f \sum_{j=1}^m \| U_j - U_{m+j} \|_{L^\infty} \nonumber \\
            &\qquad + \varphi_2 C_{\kappa} \int_0^s \bigg| P_i^{(1)}(\theta) - P_i^{(2)}(\theta) \bigg| d\theta, \label{eq:setup_for_gronwall}
\end{align}
Applying Gronwall's inequality again yields the desired result.
\end{proof}

Notice that Lemma \ref{lemma:predictor-continuity} explicitly depends on the underlying size of the spaces of the input functions which is as expected in operator approximation continuity bounds (cf. \cite[Lemma 1]{Luke_gain:2024}, \cite[Lemma 2]{Max_gain_schedule:2025}). We are now ready to present our class of approximate predictors via neural operators. 
% \subsection{Fixed point based approximate predictors} \label{subsec:iasson-approximators}
% To date, the most extensive work on numerical approximate predictors was developed in the book by Karafyllis \cite{Predictor_Feedback_For_Delay_Systems}. They proposed the following two-step iteration scheme:

% \begin{enumerate}
%     \item \textbf{Successive approximations:} Let $h> 0$ be the step size and $P_0(X_0, U, t) = X_0$ be given as an initial condition on $t \in [0, h]$. Then, consider the fixed iteration formulation scheme given by 
%     \begin{align}
%         z_{k+1}(X_0, U_1, ..., U_m, t) = X_0 + \int_0^t f(z_k(X_0, U_1, ..., U_m, s), U(s))ds, \quad k \geq 0, t \in [0, h]\,. 
%     \end{align}
%     However, notice the integral in the fixed point iteration still needs to be approximated, thus we have
%     \item \textbf{Integral discretization:}
%     For any interval $[0, \tau]$, consider $N$ subintervals of length $h = \tau/N$ and consider the interval nodes given by $ih (i =[N])$. Then the approximation for each integer $l \geq 1$ is given by 
%     \begin{align}
%         P_{i+1} = P_l(h, U(ih + s), z_i)\,, \quad i \in 0, ..., N-1, \quad P_0 = X_0\,. 
%     \end{align}
% \end{enumerate}

% Due to assumption \ref{assumption:lipschitz}, we directly obtain the following result as a consequence of \cite[Proposition 4.1]{Predictor_Feedback_For_Delay_Systems}

\subsection{Neural operator based approximate predictors} \label{subsec:neural-op-approximators}
Neural operators are neural networks defined by a specific kernel structure such that they have universal operator approximation guarantees (See \cite{Neural_Operator:2023}). There are a variety of different architectures including Deep Operator Networks (DeepONet) \cite{deeponet:2021}, Fourier Neural Operator (FNO) \cite{FNO:2021}, and even single-head attention transformers \cite{Neural_Operator:2023}. Formally, they are defined as 
\begin{definition}[Neural Operators] \label{definition:neural-operator} \cite[Section 1.2]{lanthaler2024nonlocalitynonlinearityimpliesuniversality} Let $\Omega_u \subset \mathbb{R}^{d_{u_1}}$, $\Omega_v \subset \mathbb{R}^{d_{v_1}}$ be bounded domains with Lipschitz boundary and let  $\mathcal{F}_c \subset C^0(\Omega_u; \mathbb{R}^c)$, $\mathcal{F}_v \subset C^0(\Omega_v; \mathbb{R}^v)$ be continuous function spaces.
    Given a channel dimension $d_c > 0$, we call any $\hat{\Psi}$ a neural operator given it satisfies the compositional form $\hat{\Psi} = \mathcal{Q} \circ \mathcal{L}_L \circ \cdots \circ \mathcal{L}_1 \circ \mathcal{R}$ where  $\mathcal{R}$ is a lifting layer, $\mathcal{L}_l, l=1,..., L$ are the hidden layers, and $\mathcal{Q}$ is a projection layer. That is, 
    $\mathcal{R}$ is given by 
    \begin{equation}
    \mathcal{R} : \mathcal{F}_c(\Omega_u; \mathbb{R}^c) \rightarrow \mathcal{F}_s(\Omega_s; \mathbb{R}^{d_c}), \quad c(x) \mapsto R(c(x), x)\,, 
\end{equation} where $\Omega_s \subset \mathbb{R}^{d_{s_1}}$, $\mathcal{F}_s(\Omega_s; \mathbb{R}^{d_c})$ is a Banach space for the hidden layers and $R: \mathbb{R}^c \times \Omega_u \rightarrow \mathbb{R}^{d_c}$ is a learnable neural network acting between finite-dimensional Euclidean spaces. For $l=1, ..., L$, each hidden layer is given by 
\begin{equation} \label{eq:generalNeuralOperator}
    (\mathcal{L}_l v)(x) := s \left( W_l v(x) + b_l + (\mathcal{K}_lv)(x)\right)\,, 
\end{equation}
where weights $W_l \in \mathbb{R}^{d_c \times d_c}$ and biases $b_l \in \mathbb{R}^{d_c}$ are learnable parameters, $s: \mathbb{R} \rightarrow \mathbb{R}$ is a smooth, infinitely differentiable activation function that acts component wise on inputs and $\mathcal{K}_l$ is the nonlocal operator given by 
\begin{equation} \label{eq:generalKernel}
    (\mathcal{K}_lv)(x) = \int_\mathcal{X} K_l(x, y) v(y) dy\,,
\end{equation}
where $K_l(x, y)$ is a kernel function containing learnable parameters. Lastly, the projection layer $\mathcal{Q}$ is given by 
\begin{equation}
    \mathcal{Q} : \mathcal{F}_s(\Omega_s; \mathbb{R}^{d_c}) \rightarrow \mathcal{F}_v(\Omega_v; \mathbb{R}^v), \quad s(x) \mapsto Q(s(x), y)\,, 
\end{equation}
where $Q$ is a finite dimensional neural network from $\mathbb{R}^{d_c} \times \Omega_v \rightarrow \mathbb{R}^v$. 
\end{definition}

Then, when considering the structure in Definition \ref{definition:neural-operator} as an approximator of operators defined in Definition \ref{def:delay_dependent_predictor_operators}, we obtain the following Theorem as a consequence of the continuity proved in Lemma \ref{lemma:predictor-continuity}. 
\begin{theorem}\label{thm:neural-op-approximate-predictor-theorem}
Let assumption \ref{assumption:lipschitz} hold and let $X \in \mathcal{X}$, $U_i \in C^1([0,1];\mathcal{U})$, $i \in [m]$, and $\varphi \in \Phi$, where $\mathcal{X} \subset \mathbb{R}^n$, $\mathcal{U} \subset \mathbb{R}$, and $\Phi \subset \mathbb{R}_+$ are bounded domains. For each $i \in [m]$, fix a compact set $K_i \subset \mathcal{X} \times C^1([0,1];\mathcal{U})^{m-i+1} \times \Phi$. Then, for any $\epsilon > 0$, there exist neural operator approximations $\hat{\mathcal{P}}_i: K_i \to C^1([0,1];\mathbb{R}^n)$ of the predictor operators $\mathcal{P}_i$ from Definition \ref{def:delay_dependent_predictor_operators} such that
\begin{align}
\sup_{(X,U_i,\ldots,U_m,\varphi) \in K_i}&   
\|\mathcal{P}_i(X,U_i,\ldots,U_m,\varphi)(s) \nonumber \\ & - \hat{\mathcal{P}}_i(X,U_i,\ldots,U_m,\varphi)(s)\|_{L^\infty} < \epsilon,
\end{align}
for all $X \in \mathcal{X}$, $U_i \in C^1([0,1);\mathcal{U})$, $i \in [m]$, and $\varphi \in \Phi$.
\end{theorem}


\section{Stability under uniform approximate predictors} \label{sec:stability}
We are now ready to begin studying the stability analysis of the resulting multi-input multi-delay system \eqref{eq:system} under uniform approximate predictors. To do so, we introduce the following three assumptions which are needed for the analysis of the \emph{exact} predictor feedback (\cite{Krstic:2010, Nonlinear_Control_Under_Nonconstant_Delay, Bekiaris-Liberis_Krstic:2016}) and henceforth are reasonable for our study under approximate predictors:
\begin{assumption}\label{assumption:strongly-forward-complete}
\sloppy The system $\dot{X} = f(X, \omega_1, \ldots, \omega_m)$ is strongly forward complete with respect to $\omega = (\omega_1, \ldots, \omega_m)^T$.
\end{assumption}
\begin{assumption}\label{assumption:iss}
\sloppy The system $\dot{X} = f(X, \omega_1 + \kappa_1(X), \ldots, \omega_m + \kappa_m(X))$ is input-to-state stable with respect to $\omega = (\omega_1, \ldots, \omega_m)^T$.
\end{assumption}
\begin{assumption}\label{assumption:strongly-forward-complete-2}
\sloppy The systems $\dot{X} = g_j(X, \omega_{j+1}, \ldots, \omega_m)$, for all $j=1,\dots,m$, with $g_j(X, \omega_{j+1}, \ldots, \omega_m)$ = $f(X, \kappa_1(X), \ldots, \kappa_j(X), \omega_{j+1}, \ldots, \omega_m)$, are strongly forward complete with respect to $\omega = (\omega_{j+1}, \ldots, \omega_m)^T$.
\end{assumption}

The justification for each assumption is in \cite{Bekiaris-Liberis_Krstic:2016}. We are now ready to present the main stability result. 
% Assumption \ref{assumption:strongly-forward-complete} guarantees that for every initial condition and every locally bounded input signal the corresponding solution is defined for all $t \geq 0$. In other words, the plant does not exhibit the finite escape time. Assumption \ref{assumption:iss} guarantees the robustness of the feedback system under the input disturbances, which will play a crucial role in determining the stability of the predictor feedback system. In addition to Assumption \ref{assumption:strongly-forward-complete}, Assumption \ref{assumption:strongly-forward-complete-2} guarantees that the system is strongly forward complete in between each two consecutive arrivals of control signals, and, therefore, further eliminates finite time blowups. 

\begin{theorem}\label{thm:main_result}
Let system \eqref{eq:system} satisfy Assumptions \ref{assumption:lipschitz},\ref{assumption:strongly-forward-complete}-\ref{assumption:strongly-forward-complete-2}, and assume that $\kappa_i$, $i \in [m]$, are stabilizing as in the Assumption \ref{assumption:delay-free-control-laws}. Then, for $\overline{\mathcal{B}} := \min \{ \overline{\mathcal{X}}, \overline{\mathcal{U}} \}$, there exist functions $\alpha_1 \in \mathcal{K}$ and $\beta \in \mathcal{KL}$ such that if $\epsilon < \epsilon^*$ where
\begin{align}
    \epsilon^{*}(\overline{\mathcal{B}}) := \alpha_1^{-1}(\overline{\mathcal{B}})/m, \label{eq:epsilon_star}
\end{align}
and the initial state is constrained to
\begin{align}
    \Gamma(0) \leq \Omega, \label{eq:init_state_restriction}
\end{align}
where
\begin{align}
    \Omega(\epsilon, \overline{\mathcal{B}}) := \alpha_{2}^{-1} \bigg( \overline{\mathcal{B}} - \alpha_{1}(m\epsilon)\bigg), \quad \alpha_{2}(r) = \beta(r, 0), \ r \geq 0, \label{eq:Omega}
\end{align}
then, the closed-loop solutions, under the controller with the neural operator-approximated predictor $U_i(t) = \kappa_i \bigg( \hat{P}_i(t) \bigg)$, $i \in [m]$, satisfy the $\overline{\mathcal{B}}$-semiglobal and $\epsilon$-practical stability estimate
\begin{align}
    \Gamma(t) \leq& \beta(\Gamma(0), t) + \alpha_1(m\epsilon), \quad \forall t \geq 0, \label{eq:main_inequality}
\end{align}
where
\begin{align}
    \Gamma(t) =& |X(t)| + \sum_{i=1}^m \|u_i(t)\|_{L^\infty}, \label{eq:Xi}
\end{align}
with the neural operator semiglobally trained with $\overline{\mathcal{B}} > \alpha_2(\Omega)$ and $\epsilon \in (0, \alpha_1^{-1}(\overline{\mathcal{B}})/m)$.
\end{theorem}

Before proving the result, we briefly discuss its implications. The result states that for all $\epsilon < \epsilon^*$, there exists a region of attraction $\Omega$ whose size grows inversely with $\epsilon$, such that all initial states $\Gamma(0) \leq \Omega$ are $\epsilon$-practically stable. Conversely, for $\epsilon > \epsilon^*$, no such region exists, and the system becomes unstable for any neural network predictor with large errors. Thus, the result has two key perspectives: for neural network engineers, achieving a sufficiently small $\epsilon$ is essential to stabilize the system; for control verification engineers, accurately estimating $\epsilon$ is crucial, as it directly determines both the region of attraction and the size of the invariant stable set. Such estimation is beyond the scope of this study but may be performed using, for example, the $L_2$ testing error.

\begin{proof}
    The proof of Theorem \ref{thm:main_result} will be split over a series of technical Lemmas across two subsections. We begin by introducing an equivalent representation of \eqref{eq:system} representing $m$ delays as individual transport PDEs.
\subsection{Transport PDE equivalent delay system}
Consider the setting in which $U_i$ is invoked with the approximated predictions as in Theorem \ref{thm:neural-op-approximate-predictor-theorem}. Thus, the approximate predictors are given by the operator approximations as 
\begin{subequations}
\begin{align}
    \hat{P}_1(t) =& \hat{\mathcal{P}}_1(X(t), T_{D_1,0}(t)U_1, ..., T_{D_m,D_{m-1}}(t)U_m, D_1)(1), \\
    \hat{P}_2(t) =& \hat{\mathcal{P}}_2(\hat{P}_1(t), T_{D_{2,1},0}(t)U_2, \nonumber \\ &..., T_{D_{m,1},D_{m,2}}(t)U_m, D_{2,1})(1), \\
    &\ \hspace{3cm} \vdots \nonumber \\
    \hat{P}_m(t) =& \hat{\mathcal{P}}_m(\hat{P}_{m-1}(t), T_{D_{m,m-1},0}(t)U_m, D_{m,m-1})(1).
\end{align}
\end{subequations}
Then,  system \eqref{eq:system} can be written equivalently in its ODE-PDE form, originally developed in \cite{Krstic:2010}, as
\begin{subequations}\label{eqs:plant_ode_pde}
    \begin{align}
        \dot{X}(t) =& f(X(t), u_1(0,t), \ldots, u_m(0,t)), \label{eq:state} \\
        \frac{\partial}{\partial t} u_i(x,t) =& \frac{\partial}{\partial x} u_i(x,t), \quad (x,t) \in [0,D_i) \times \mathbb{R}_+,\label{eq:delay_pde} \\
        u_i(D_i,t) =&\kappa_i(\hat{P}_i(t)), \quad i \in [m]. \label{eq:pde_boundary}
    \end{align}
\end{subequations}
The rewritten system is an ODE coupled with $m$ different transport PDEs, where the analytical solution of each transport PDE is given as
\begin{align}
    u_i(x,t) = U_i(t+x-D_i), \quad x \in [0, D_i),
\end{align}
$i \in [m]$. Notice that the transport PDEs \eqref{eq:delay_pde} are defined on growing spatial domains as the number of inputs increases. Hence, the time for the input signal to reach from the input boundary at position $D_i$ to the opposite side at $0$ is exactly $D_i$ seconds with the unit transport speed. 






The corresponding predictors then become, in PDE notation, 
\begin{subequations}\label{eqs:ps}
    \begin{align}
        p_1(x,t) =& X(t) + \int_0^x f \bigg( p_1(y,t), u_1(y,t), \ldots, u_m(y,t) \bigg) dy, \label{eq:p1(x,t)}  \nonumber \\
            &\quad x \in [0, D_1) \\
        p_2(x,t) =& p_1(D_1,t) + \int_{D_1}^x f \bigg( p_2(y,t), \kappa_1(p_2(y,t)), u_2(y,t),\nonumber \\ & \ldots, u_m(y,t) \bigg) dy, \label{eq:p2(x,t)} \quad x \in [D_1, D_2) \\
        &\hspace{2cm} \vdots \nonumber \\
        p_m(x,t) =& p_{m-1}(D_{m-1},t) + \int_{D_{m-1}}^x f \bigg( p_m(y,t), \nonumber \\ &  \kappa_1(p_m(y,t)), \ldots, \kappa_{m-1}(p_m(y,t)), u_m(y,t) \bigg) dy, \label{eq:pm(x,t)} \nonumber \\
            &\quad x \in [D_{m-1}, D_m) .
    \end{align}
\end{subequations}
Notice that, the predictors $p_1(D_1, t), ..., p_m(D_m, t)$ are equivalent to the original predictors $P_1(t), ..., P_m(t)$ in \eqref{eqs:P_system}. The key advantage of rewriting \eqref{eq:system} in the form of \eqref{eq:state}, \eqref{eq:delay_pde}, \eqref{eq:pde_boundary} is that analyzing the ODE-PDE cascade becomes significantly simpler then the original ODE under delayed inputs.

\subsection{Stability of the ODE-PDE system}
To analyze the system, we begin by introducing the following backstepping transformation as in \cite{Bekiaris-Liberis_Krstic:2016} to transform the ODE-PDE cascade into what we call the ``target system". The key advantage of the target system is that it is feasible to show that the system is ISS-stable \cite{iasson-iss}. This stability result for the target system can then be transformed into a stability result for the original system via an inverse backstepping transformation, completing the main result. 
\begin{lemma}\label{lemma:backstepping-transform-to-target-system}
The backstepping transformations for the transport PDEs $u_i$ defined by
\begin{subequations}
\begin{align}
    w_1(x,t) =& u_1(x,t) - \kappa_1\bigg(p_1(x,t)\bigg), \quad x \in [0, D_1) \label{eq:w1} \\
    w_2(x,t) =& u_2(x,t) - 
        \begin{cases} 
            \kappa_2\bigg(p_1(x,t)\bigg), \quad x \in [0, D_1) \\
            \kappa_2\bigg(p_2(x,t)\bigg), \quad x \in [D_1, D_2)
        \end{cases} \label{eq:w2} \\
    &\ \vdots \nonumber \\
    w_m(x,t) =& u_m(x,t) - 
        \begin{cases} 
            \kappa_m\bigg(p_1(x,t)\bigg), \quad x \in [0, D_1) \\
            \kappa_m\bigg(p_2(x,t)\bigg), \quad x \in [D_1, D_2) \\
            \ \vdots \\
            \kappa_m\bigg(p_m(x,t)\bigg), \quad x \in [D_{m-1}, D_m)
        \end{cases}, \label{eq:wm}
\end{align}
\end{subequations}
transform the plant \eqref{eqs:plant_ode_pde} to the target system $(i \in [m])$
\begin{subequations}\label{eqs:target}
    \begin{align}
        \dot{X}(t) &= f \bigg( X(t), \kappa_1(X(t)) + w_1(0,t),\nonumber \\ & \dots, \kappa_m(X(t)) + w_m(0,t) \bigg),  \label{eq:target_ode} \\
        \partial_t w_i(x,t) &= \partial_x w_i(x,t), \quad (x,t) \in (0,D_i) \times [0,\infty) \label{eq:target_pde} \\
        w_i(D_i,t) &= \kappa_i(\hat{P}_i(t)) - \kappa_i(P_i(t)), \quad t \in [0,\infty). \label{eq:target_pde_boundary}
    \end{align}
\end{subequations}
\end{lemma}

\begin{proof}
    The proof follows exactly that of \cite[Lemma 1]{Bekiaris-Liberis_Krstic:2016} except that the boundary condition on \eqref{eq:pde_boundary} is with the approximate predictor and henceforth the resulting boundary condition of the target system is obtained via direct substitution, yielding \eqref{eq:target_pde_boundary}. 
\end{proof}


Additionally, there exists an invertible set of transformations from the $w$ to the $u$ system characterized by

\begin{lemma}\label{lemma:inverse-backstepping-transforms}
    For all $i \in [m]$, the inverse backstepping transformation for the PDEs $w_i$ defined by 
    \begin{subequations}
    \begin{align}
        u_1(x,t) =& w_1(x,t) + \kappa_1\bigg(\pi_1(x,t)\bigg), \quad x \in [0, D_1) \label{eq:inv_w1} \\
        u_2(x,t) =& w_2(x,t) +
            \begin{cases}
                \kappa_2\bigg(\pi_1(x,t)\bigg), \quad x \in [0, D_1) \\
                \kappa_2\bigg(\pi_2(x,t)\bigg), \quad x \in [D_1, D_2)
            \end{cases} \label{eq:inv_w2} \\
        &\ \vdots \nonumber \\
        u_m(x,t) =& w_m(x,t) -
            \begin{cases}
                \kappa_m\bigg(\pi_1(x,t)\bigg), \quad x \in [0, D_1) \\
                \kappa_m\bigg(\pi_2(x,t)\bigg), \quad x \in [D_1, D_2) \\
                \ \vdots \\
                \kappa_m\bigg(\pi_m(x,t)\bigg),\\ \qquad x \in [D_{m-1}, D_m)
            \end{cases} \label{eq:inv_wm}
    \end{align}
    \end{subequations}
    where
    \begin{subequations}
    \begin{align}
        \pi_1(x,t) =& X(t) + \int_0^x f \bigg (\pi_1(y,t), w_1(y,t) + \kappa_1 \bigg( \pi_1(y,t) \bigg), \nonumber \\ & \ldots, w_m(y,t) + \kappa_m \bigg( \pi_1(y,t) \bigg) \bigg) dy, \label{eq:pi1} \\
            &\qquad x \in [0, D_1] \nonumber \\
        \pi_i(x,t) =& \pi_{i-1}(D_{i-1},t) + \int_{D_{i-1}}^x f \bigg( \pi_i(y,t), \nonumber \\ & w_1(y,t) + \kappa_1 \bigg( \pi_i(y,t) \bigg), \ldots, \nonumber \\ &  w_m(y,t) + \kappa_m \bigg( \pi_m(y,t) \bigg) \bigg) dy, \label{eq:pii} \\
            &\qquad x \in [D_{i-1}, D_i], \quad i=2,\dots,m. \nonumber
    \end{align}
    \end{subequations}
    transforms the plant \eqref{eq:target_ode}, \eqref{eq:target_pde}, \eqref{eq:target_pde_boundary} into the plant \eqref{eq:state}, \eqref{eq:delay_pde}, \eqref{eq:pde_boundary}. 
\end{lemma}

\begin{proof}
    As in Lemma \ref{lemma:backstepping-transform-to-target-system}, the proof follows \cite[Lemma 2]{Bekiaris-Liberis_Krstic:2016} with the appropriate substitutions of the boundary condition due to the approximate predictors. 
\end{proof}


We are now ready to provide an ISS stability result for the target ODE-PDE system. 


\begin{lemma}\label{lem:lemma_3_of_6}
Let the system \eqref{eqs:target} satisfy the Assumption \ref{assumption:iss}. Then, there exists a class $\mathcal{KL}$ function $\beta_1$ and a class $\mathcal{K}$ function $\gamma_1$ such that the following holds for all $t \geq 0$
\begin{align}
    \breve{\Gamma}(t) \leq& \nonumber  \beta_1 \bigg( \breve{\Gamma}(0), t \bigg) \\ &+ \gamma_1 \bigg( \sup_{0 \leq \tau \leq t} \bigg\{ \sum_{i=1}^m |\kappa_i(\hat{P}_i(t)) - \kappa_i(P_i(t))| \bigg\} \bigg), \label{eq:lemma3_Gamma_bar}
\end{align}
where
\begin{align}
    \breve{\Gamma}(t) :=& |X(t)| + \sum_{i=1}^m \|w_i(t)\|_{L^\infty}. \label{eq:Gamma_bar_def}
\end{align}
\end{lemma}
\begin{proof}
From Assumption \ref{assumption:iss} and the definition of input-to-state stability in \cite{SontagISS:1995}, it follows that there exist function $\beta_2 \in \mathcal{KL}$ and $\gamma_2 \in \mathcal{K}$ such that
\begin{align}
    |X(t)| \leq& \beta_2(|X(0)|, t) + \gamma_2 \bigg( \sup_{0 \leq \tau \leq t} \bigg\{ \sum_{i=1}^m |w_i(0,\tau)| \bigg\} \bigg) \label{eq:w_iss}\,, 
\end{align}
for all $t \geq 0$. Using the ISS estimate of the transport PDE from \cite[Lemma 6]{Luke_Peijia:2025}, for any constant $c> 0$, for all $i \in [m]$, and for all times $t > 0$, we have that 
\begin{align}
    \|w_i(t)\|_{L^\infty} \leq& \exp(c(D_i-t)) \|w_i(0)\|_{L^\infty} \nonumber  \\ &+ \exp(cD_i) \sup_{0 \leq \tau \leq t} |w_i(D_i,\tau)|. \label{eq:wi_pde_iss}
\end{align}
From \eqref{eq:w_iss}, we have
\begin{align}
    |X(t)| \leq& \beta_2(|X(0)|, t) + \gamma_2 \bigg( \sup_{0 \leq \tau \leq t} \bigg\{ \sum_{i=1}^m |w_i(0,\tau)| \bigg\} \bigg) \nonumber \\
        \leq& \beta_2(|X(0)|, t) + \gamma_2 \bigg( \sup_{0 \leq \tau \leq t} \bigg\{ \sum_{i=1}^m \|w_i(\tau)\|_{L^\infty} \nonumber \\ & + \sum_{i=1}^m |w_i(D_i,\tau)| \bigg\} \bigg) \label{eq:lemma3_x_iss},
\end{align}
and, from \eqref{eq:wi_pde_iss} applied to all $i \in [m]$, we have
\begin{align}
    \sum_{i=1}^m \|w_i(t)\|_{L^\infty}  %\leq& \sum_{i=1}^m \exp(c(D_i-t)) \|w_i(0)\|_{L^\infty}  \nonumber \\ & + \sum_{i=1}^m \exp(cD_i) \sup_{0 \leq \tau \leq t} \bigg|w_i(D_i,\tau)\bigg| \nonumber \\
        \leq& \beta_3 \bigg( \sum_{i=1}^m \|w_i(0)\|_{L^\infty}, t \bigg) \nonumber \\ & + \gamma_3 \bigg( \sup_{0 \leq \tau \leq t} \bigg\{ \sum_{i=1}^m |w_i(D_i,\tau)| \bigg\} \bigg), \label{eq:lemma3_w_iss}
\end{align}
where $\beta_3 \in \mathcal{KL}$, and $\gamma_3 \in \mathcal{K}$. Now, treating $X(t)$ and $\sum_{i=1}^m \|w_i(t)\|_{L^\infty}$ as two interconnected subsystems with $\sup_{0 \leq \tau \leq t} \sum_{i=1}^m|w_i(D_i,\tau)|$ as their input, we can apply the classical results of cascaded ISS systems (See \cite[Lemme C.4]{KKK} or \cite{402246}) using the forms \eqref{eq:lemma3_x_iss}, \eqref{eq:lemma3_w_iss} to obtain 
% Now, let $(a, b)$ be the column concatenation of the vectors $a, b$. Then, by reproducing the argument from Lemma C.4 in \cite{KKK}, we obtain
% \begin{align}
%     \bigg| \bigg(X(t)&, \sum_{i=1}^m \|w_i(t)\|_{L^\infty} \bigg) \bigg|\nonumber \\  \leq& \beta_4 \bigg( \bigg| \bigg( X(0), \sum_{i=1}^m \|w_i(0)\|_{L^\infty} \bigg) \bigg|, t \bigg) \nonumber \\
%         &\qquad + \gamma_4 \bigg( \sup_{0 \leq \tau \leq t} \bigg\{ \sum_{i=1}^m |w_i(D_i,\tau)| \bigg\} \bigg), \label{eq:after_c4}
% \end{align}
% where $\beta_4 \in \mathcal{KL}$, and $\gamma_4 \in \mathcal{K}$. Furthermore, using the facts that $|x_1| \leq |(x_1,x_2)|$, $|x_2| \leq |(x_1,x_2)|$, and $|(x_1,x_2)| \leq |x_1| + |x_2|$, we obtain the following two inequalities
% \begin{align}
%     |X(t)| \leq& \beta_4 \bigg( |X(0)| + \sum_{i=1}^m \|w_i(0)\|_{L^\infty}, t \bigg)  \nonumber \\ &+ \gamma_4 \bigg( \sup_{0 \leq \tau \leq t} \bigg\{ \sum_{i=1}^m |w_i(D_i,\tau)| \bigg\} \bigg), \label{eq:x_after_c4}
% \end{align}
% and
% \begin{align}
%     \sum_{i=1}^m \|w_i(0)\|_{L^\infty} \leq& \beta_4 \bigg( |X(0)| + \sum_{i=1}^m \|w_i(0)\|_{L^\infty}, t \bigg) \nonumber \\ & + \gamma_4 \bigg( \sup_{0 \leq \tau \leq t} \bigg\{ \sum_{i=1}^m |w_i(D_i,\tau)| \bigg\} \bigg). \label{eq:w_after_c4}
% \end{align}
% It follows from \eqref{eq:x_after_c4} and \eqref{eq:w_after_c4} that
\begin{align}
    |X(t)|\nonumber  + \sum_{i=1}^m \|w_i(0)\|_{L^\infty} \leq& \beta_1 \bigg( |X(0)| + \sum_{i=1}^m \|w_i(0)\|_{L^\infty}, t \bigg) \nonumber \\&\!\!\!+ \gamma_1 \bigg( \sup_{0 \leq \tau \leq t} \bigg\{ \sum_{i=1}^m |w_i(D_i,\tau)| \bigg\} \bigg), \label{eq:x_and_w_after_c4}
\end{align}
where $\beta_1 \in \mathcal{KL}$, and $\gamma_1 \in \mathcal{K}$. Lastly, by defining the function $\breve{\Gamma}(t) = |X(t)| + \sum_{i=1}^m \|w_i(t)\|_{L^\infty}$ and using \eqref{eq:target_pde_boundary} completes the proof.
\end{proof}

To complete the proof of Theorem \ref{thm:main_result}, the exact same approach to convert from the bound on the target system to the original system using \cite[Lemmas 4, 5, 6]{Bekiaris-Liberis_Krstic:2016} can be applied to \eqref{eq:lemma3_Gamma_bar} using properties of class $\mathcal{K}$ functions. Thus, there exists $\beta_5 \in \mathcal{KL}$, $\gamma_5 \in \mathcal{KL}$ such that
\begin{align}
    \Gamma(t) \leq& \beta_5 \bigg( \Gamma(0), t \bigg) \nonumber  \\ & + \gamma_5 \bigg( \sup_{0 \leq \tau \leq t} \bigg\{ \sum_{i=1}^m |\kappa_i(\hat{P}_i(t)) - \kappa_i(P_i(t))| \bigg\} \bigg), \label{eq:original-system-bound}
\end{align}
Using the Lipschitz constant $C_\kappa$ as in \eqref{eq:c_kappa} with the universal approximation theorem (Theorem \ref{thm:neural-op-approximate-predictor-theorem}), and property of class $\mathcal{K}$ functions, one achieves Theorem \ref{thm:main_result}. 
\end{proof}


\section{Numerical simulations} \label{sec:numerical}
To validate our theoretical results, we develop neural operator approximate predictors for the unicycle system with distinct steering and drive delays. Consider the following: 
\begin{subequations}
\begin{align}
    x_(t) &= U_2(t-D_2)\cos(\theta(t)) \label{eq:unicycle_x1_dot} \\
    y(t) &= U_2(t-D_2)\sin(\theta(t)) \label{eq:unicycle_x2_dot} \\
    \theta(t) &= U_1(t-D_1), \label{eq:unicycle_x3_dot}
\end{align}
\end{subequations}
where $x, y$ are the 2D coordinate positions of the unicycle, $\theta$ is spatial orientation, $U_1$ is the turning rate, $U_2$ is linear speed, and $D_2$ is the longer breaking delay and $D_1$ is the short steering delay. We consider the time-varying controller of \cite{POMET1992147} given as presented in \cite[Section V]{Bekiaris-Liberis_Krstic:2016}.

In our numerical experiment, we consider delays of $D_1 = 0.25$s and $D_2= 0.6$s with a time step $dt=0.001$s from $0$ to $T=10$s of simulation time. 
As both a baseline and to create a dataset for training the neural operator, we simulate $100$ trajectories using the numerical approach of \cite{Predictor_Feedback_For_Delay_Systems}, applying additive uniform noise $\mathcal{U}(-0.2, 0.2)$ to the state at each timestep. From each trajectory, we collect all predictor input/output pairs, resulting in a total dataset of 200,000 pairs. All trajectories start from the same initial conditions $(x, y, \theta) = \bm{0}$, with zero initial input $U_j(s) = 0$, $s \in [-D_j, 0]$, $j \in [2]$. The dataset generation process takes about an hour on a standard laptop. For training, we consider two popular neural operator architectures—FNO and DeepONet—which are trained in a supervised manner(10 minutes, Nvidia 4060 GPU). \textcolor{blue}{See \cite{ProjectGitHub} for all code, datasets, models, and hyperparameters.}

% \subsection{Dataset construction and model training}
% Before proceeding to describe the database construction process, it is necessary to briefly revisit Theorem ({\color{red} To be referenced}), which constitutes the core part of our work concerning numerical approximations via machine learning (ML) models. Specifically, the theorem guarantees uniform accuracy only for models that satisfy the universal approximation theorem ({\color{red} Appendix ????}). Therefore, we will analyze two types of ML models that have the required property: DeepONet \cite{deeponet:2021} and FNO (Fourier Neural Operator) \cite{FNO:2021}.

% We start by creating the database using numerical simulations based on the Fixed point iteration method, which provides the ground truth for training and validation. First, we simulate 100 trajectories, each spanning $T = 10\,\text{s}$ with a time step of $\Delta t = 0.001$ s. At each time step, we recompute the prediction horizon for all states using a spatial resolution of $\Delta x = 0.01$. Each trajectory starts in the same initial condition $(x, y, \theta) = [1, 1, 1]$ with uniform noise between $-0.2$ and $0.2$ to ensure the uniqueness of all trajectories. Furthermore, we assume that the initial control input history is zero, that is, $U_i(t) = 0$, $t \in [-D_i,0)$ for all $i \in [m]$.

% To train the neural operators, we performed a hyperparameter search for the optimization parameters (epochs, learning rate, weight decay, etc.) while using AdamW \cite{AdamW} as our optimizer with an exponential learning rate scheduler, as is standard practice for neural operators. We train using the standard $L_2$ loss and refer the reader to the Jupyter notebook examples on GitHub \cite{ProjectGitHub} for full parameter details.

% Before continuing, it is important to emphasize that system \eqref{eqs:plant_ode_pde} can be represented in an appropriate PDE formulation, where all control signals are scaled to the interval $[0, 1]$. This scaling ensures that we only need to train a single model, which can then be used for all $m$ predictors. The only adjustment required is an appropriate modification of the length of the prediction horizon.

% \begin{figure*}[t]
%     \centering
%     \includegraphics[trim=0 180 0 25, clip, width=\linewidth]{sections/section3/plots/single_trajectory_all.png}
%     \caption{Example evaluation trajectories with fixed point, DeepONet, and FNO approximated predictors. \todo{legend needed} }
%     \label{fig:single_trajectory}
%     \vspace{-1em}
% \end{figure*}

\begin{figure*}[t]
    \centering
    \includegraphics[trim=0 5px 0 7px, clip]{single_trajectory_all_upper_half.png}
    \caption{Example evaluation trajectories with fixed point, DeepONet, and FNO approximated predictors.}
    \label{fig:single_trajectory}
    \vspace{-1em}
\end{figure*}

% \subsection{Results}
% We choose $D_1 = 0.25$ and $D_2 = 0.6$ for our demonstrative example.
In Table \ref{tab:performance_errors}, we present the performance of the neural approximated predictors.
%from Definition \ref{def:delay_dependent_predictor_operators} throughout both training and validation datasets corresponding to time step $\Delta t=0.001$ and spatial resolution $\Delta x=0.01$. 
All approaches achieve strong performance with errors on the magnitude of $\sim\!10^{-6}$. 
Furthermore, to accurately evaluate the neural approximated predictor operator in the feedback loop, we simulate 25 new trajectories for each approximation approach and present the average trajectory residuals with respect to the stationary equilibrium point, $(x,y,\theta)=\bm{0}$ in Table \ref{tab:performance_errors}. Furthermore, to showcase the similarity between both approaches, we provide an example of one of the $25$ trajectories in Figure \ref{fig:single_trajectory}.

Table \ref{tab:performance_times} highlights the computational advantages of neural operator approximation of predictors over fixed point iteration, particularly at finer spatial resolutions. While fixed point iteration’s computation time increases significantly as resolution improves, FNO maintains nearly constant runtimes, and DeepONet shows moderate variation. These traits enable both neural operators to outperform traditional methods as resolution becomes finer. Notably, their speedups are largely due to deployment on specialized GPU hardware compared to CPU-based fixed point iteration.
% and presente in column four of Table \ref{tab:performance_errors}. Additionally, all initial conditions for those testing trajectories are sampled from the same distribution used for generating training datasets. Here, it is important to emphasize that these neural predictors are trained offline and deployed online, in the sense that their inference constitutes the crucial component of the predictor-feedback controller.

% Table \ref{tab:performance_times} highlights the computational advantages of FNO and DeepONet over the Fixed point iteration, particularly for finer spatial resolutions. As they decrease, the computational gain of these methods becomes pronounced. Fixed point iteration exhibits a linear scaling in computation time with respect to $\Delta x$, as seen in the increasing values from 8.762 ms ($\Delta x=0.01$) to 218.740 ms ($\Delta x=0.0005$). In contrast, FNO's computation time remains nearly constant across all resolutions (e.g., ~31 ms), requiring the same computational effort regardless of temporal resolution. On the other hand, DeepONet shows non-linear scaling, with computation times fluctuating between 7.245 ms and 23.126 ms. This behavior enables FNO and DeepONet to achieve computational advantages for finer spatial resolutions, where traditional method such as Fixed point iteration becomes increasingly expensive. To emphasize, the performance gains of FNO and DeepONet stem from their deployment on specialized hardware (GPUs), which achieves significant speedups compared to traditional CPU-based solutions. 
% Additionally, the simulation examples of all three methods, starting from the same/distinct initial conditions, are presented in Figure \ref{fig:single_trajectory} and Figure \ref{fig:multiple_trajectory}, respectively.

% In summary, the numerical simulations demonstrate that both FNO and DeepONet achieve high accuracy (errors $\sim\!10^{-6}$) in approximating the predictor operators while offering significant computational advantages over traditional Fixed point iteration method, particularly for fine spatial resolutions. The consistent performance across training, validation, and dynamical testing trajectories confirms the robustness of these neural operators. Their ability to maintain low computational costs regardless of temporal resolution (unlike Fixed point iteration) highlights their suitability for real-time control applications with delays. Further details on implementation and extended results are provided on GitHub \cite{ProjectGitHub}.

\begin{table}[H]
\centering
\caption{
  The performance of predictor approximation approaches across training and validation datasets. The average trajectory tracking error is the summed $L^2$ error over time averaged over the 25 different test scenarios (with randomized initial conditions).}
  %, as well as the average trajectory $L_2$ error across the testing dataset, with constant and distinct delays for unicycle stabilization with time step $\Delta t=0.001$ and $\Delta x=0.01$.
\label{tab:performance_errors}
\resizebox{0.5\textwidth}{!}{
\begin{tabular}{@{}*{5}{c}@{}}
    \toprule
    \multirow{2}{*}{\makecell[c]{Approximation\\methods}} & 
    \multirow{2}{*}{\makecell[c]{Parameters}} & 
    \multirow{2}{*}{\makecell[c]{Training $L_2$\\error $(\times10^{-6})\downarrow$}} & 
    \multirow{2}{*}{\makecell[c]{Validation $L_2$\\error $(\times10^{-6})\downarrow$}} & 
    \multirow{2}{*}{\makecell[c]{Avg. trajectory\\$L_2$ error $\downarrow$}} \\
    & & & & \\
    \midrule
    \makecell[c]{Fixed point\\approximation} & \makecell[c]{---} & \makecell[c]{---} & \makecell[c]{---} & 5.131 \\
    FNO & 42211 & 0.33 & 0.25 & 5.144 \\
    DeepONet & 1124824 & 3.00 & 2.46 & 4.874 \\
    \bottomrule
\end{tabular}
}
\end{table}

\vspace{-1em}

\begin{table}[H]
\centering
\caption{
  Computation time (ms) averaged over $1,000$ predictions. % for various step sizes and speedups relative to the Fixed point approximation.
}
\label{tab:performance_times}
\resizebox{0.5\textwidth}{!}{
\begin{tabular}{@{}*{6}{c}@{}}
    \toprule
    \multirow{2}{*}{\makecell[c]{Approximation\\methods}} & 
    \multicolumn{4}{c}{\makecell[c]{Computation time (ms)}} & 
    \multirow{2}{*}{\makecell[c]{Speedup $\uparrow$}} \\
    \cmidrule(lr){2-5}
    & \makecell[c]{dx = 0.01} & \makecell[c]{dx = 0.005} & \makecell[c]{dx = 0.001} & \makecell[c]{dx = 0.0005} & \\
    \midrule
    \makecell[c]{Fixed point\\approximation} & 8.762 & 17.735 & 97.676 & 218.740 & \makecell[c]{---} \\
    FNO & 30.850 & 29.664 & 34.653 & 31.093 & 7.04$\times$ \\
    DeepONet & 7.245 & 9.704 & 10.816 & 23.126 & 9.46$\times$ \\
    \bottomrule
\end{tabular}
}
\end{table}

% In summary, we briefly discuss the hurdles of designing a neural approximated predictor operator with distinct constant delays as well as the advantages of this approach. First, we observe that as the gap between delays increases, the amount of data necessary to achieve the same approximation accuracy increases. This is to be expected since the neural network has to learn multiple dynamics that correspond to the period between two consecutive control hits. On the other hand, the advantage of neural operator predictors over the traditional successive approximations approach becomes more pronounced in more challenging settings. One such example is presented in \cite{Luke_Peijia:2025}, where the cost of executing manipulator dynamics is significantly greater than the cost of a single forward pass in a neural network. This presents the most important trade-off in applying neural networks to control problems. One will have to spend more time on development, and it will be paid off through computational efficiency. In the end, such compromises are naturally unavoidable when implementing real-time controllers in complex systems.

% \begin{table*}[htbp]
% \centering
% \caption{
%   The performance of predictor approximation approaches across training and validation datasets, as well as the average trajectory $L_2$ error across the testing dataset, with constant and distinct delays for unicycle stabilization with time step $\Delta t=0.01$ (left) and computation times for various step sizes (right).
% }
% \label{tab:performance}
% \setlength{\tabcolsep}{4pt}
% \resizebox{\textwidth}{!}{
% \begin{tabular}{@{}*{9}{c}@{}}
%     \toprule
%     \multirow{2}{*}{\makecell[c]{Approximation\\Method}} & 
%     \multirow{2}{*}{\makecell[c]{Parameters}} & 
%     \multirow{2}{*}{\makecell[c]{Training $L_2$\\error $(\times10^{-6})\downarrow$}} & 
%     \multirow{2}{*}{\makecell[c]{Validation $L_2$\\error $(\times10^{-6})\downarrow$}} & 
%     \multirow{2}{*}{\makecell[c]{Avg. trajectory\\$L_2$ error $\downarrow$}} & 
%     \multicolumn{4}{c}{\makecell[c]{Computation time (ms)}} \\
%     \cmidrule(lr){6-9}
%     & & & & & \makecell[c]{dx = 0.01} & \makecell[c]{dx = 0.005} & \makecell[c]{dx = 0.001} & \makecell[c]{dx = 0.0005} \\
%     \midrule
%     \makecell[c]{Successive\\approximations} & \makecell[c]{---} & \makecell[c]{---} & \makecell[c]{---} & 4.633 & 8.762 & 17.735 & 97.676 & 218.740 \\
%     FNO & 42211 & 0.33 & 0.25 & 5.121 & 30.850 & 29.664 & 34.653 & 31.093 \\
%     DeepONet & 1124824 & 3.00 & 2.46 & 4.420 & 7.245 & 9.704 & 10.816 & 23.126 \\
%     \bottomrule
% \end{tabular}
% }
% \end{table*}





% \begin{figure*}[htbp]
%     \centering
%     \includegraphics{sections/section3/plots/multiple_trajectories_all.png}
%     \caption{Example of multiple trajectories of the unicycle system with the Fixed point, DeepONet, and FNO approximated predictors, starting from the distinct initial conditions.}
%     \label{fig:multiple_trajectory}
% \end{figure*}



% For a more detailed explanation of the recursive construction of these predictors, see the work of \cite{Bekiaris-Liberis_Krstic:2016}. In addition, to proceed further, we need the following assumptions
% \begin{assumption}\label{ass:Assumption1}
% The system $\dot{X} = f(X, \omega_1, \ldots, \omega_m)$ is strongly forward complete with respect to $\omega = (\omega_1, \ldots, \omega_m)^T$.
% \end{assumption}
% \begin{assumption}\label{ass:Assumption2}
% The system $\dot{X} = f(X, \omega_1 + \kappa_1(X), \ldots, \omega_m + \kappa_m(X))$ is input-to-state stable with respect to $\omega = (\omega_1, \ldots, \omega_m)^T$.
% \end{assumption}
% \begin{assumption}\label{ass:Assumption3}
% The systems $\dot{X} = g_j(X, \omega_{j+1}, \ldots, \omega_m)$, for all $i=1,\dots,m$, with $g_j(X, \omega_{j+1}, \ldots, \omega_m)$ = $f(X, \kappa_1(X), \ldots, \kappa_j(X), \omega_{j+1}, \ldots, \omega_m)$, are strongly forward complete with respect to $\omega = (\omega_{j+1}, \ldots, \omega_m)^T$.
% \end{assumption}
% Those three assumptions are necessary for analyzing systems with long input delays (\cite{Krstic:2010}, \cite{Nonlinear_Control_Under_Nonconstant_Delay}, and \cite{Bekiaris-Liberis_Krstic:2016}). The Assumption \ref{ass:Assumption1} guarantees that for every initial condition and every locally bounded input signal the corresponding solution is defined for all $t \geq 0$, in other words, the plant does not exhibit the finite escape time. The Assumption \ref{ass:Assumption2} guarantees the robustness of the feedback system under the input disturbances, which will play a crucial role in determining the stability of the predictor feedback system. In contrast to the \cite{Bekiaris-Liberis_Krstic:2016}, this assumption cannot be relaxed to global asymptotic stability since the disturbance can destabilize the system even after all delays have been compensated. In addition to the Assumption \ref{ass:Assumption1}, the Assumption \ref{ass:Assumption3} guarantees that the system is strongly forward complete in between each two consecutive arrivals of control signals, and, therefore, ensuring no finite escape time.
%  Also, it can be seen that the functions $p_i$, $i \in [m]$, satisfy the following ODEs in $x$
% \begin{align}
%     \partial_xp_1(x,t) =& f \bigg( p_1(y,t), u_1(y,t), \ldots, u_m(y,t) \bigg), \quad x \in [0, D_1) \\
%     \partial_xp_2(x,t) =& f \bigg( p_2(y,t), \kappa_1(p_2(y,t)), u_2(y,t), \ldots, u_m(y,t) \bigg), \quad x \in [D_1, D_2) \\
%     &\ \vdots \nonumber \\
%     \partial_xp_m(x,t) =& f \bigg( p_m(y,t), \kappa_1(p_m(y,t)), \ldots, \kappa_{m-1}(p_m(y,t)), u_m(y,t) \bigg), \quad x \in [D_{m-1}, D_m),
% \end{align}
% with initial conditions
% \begin{align}
%     p_1(0,t) =& X(t) \label{eq:p1(0,t)} \\
%     p_2(D_1,t) =& p_1(D_1,t) \label{eq:p2(0,t)} \\
%     &\ \vdots \nonumber \\
%     p_m(D_{m-1},t) =& p_{m-1}(D_{m-1},t) \label{eq:pm(0,t)}.
% \end{align}
% Note that with this representation of predictors, we have that
% \begin{equation}
%     U_i(t) = \kappa_i(p_i(D_i,t)), \quad i=1,\dots,m \label{eq:control_laws}.
% \end{equation}
% The controllers in \eqref{eq:control_laws} provide global asymptotic stability of the system \eqref{eq:system}. For a detailed explanation, we refer the reader to Theorem 1 in \cite{Bekiaris-Liberis_Krstic:2016}, which is restated in the Appendix.

% \section{Stability under NeuralOP Approximated Predictor Feedback}
% The first neural operators have been introduced in \cite{Neural_Operator:2023}, and, since then, they have gained popularity in many areas of science and engineering. It was \cite{Luke_Peijia:2025} in which neural operators were first utilized for predictor-feedback control of systems with constant input delay. We extend their work and generalize to multi-input systems with distinct input delays. When it comes to the organization of this section, we start by defining multiple neural approximated predictor operators. Then, we prove that each defined operator is continuous under the assumption of Lipschitz dynamics of the function $f$. Lastly, the continuity of the operators enables us to invoke the universal operator approximation theorem from \cite{LLS2024}, which provides us with the necessary approximation guarantees.

% \subsection{Neural Approximated Predictor Operators}
% The prediction horizon in multi-input systems with distinct delays is naturally divided into multiple intervals, each corresponding to the interval between two consecutive delays $D_{i-1}$ and $D_i$ as illustrated in equations \eqref{eqs:ps}. Crucially, each interval exhibits its own distinct dynamics due to the varying combinations of control inputs and feedback terms. To address this complexity, the most practical approach is to train $m$ separate neural operators, each dedicated to predicting the state within its respective interval.
% \begin{definition}\label{def:delay_dependent_predictor_operators}
% (\textbf{Delay dependent predictor operators for multi-input systems}) Let $Q \in \mathbb{R}^n$, and for each $i \in [m]$, let $U_i \in C^1([0,1];\mathbb{R})$. Let $\varphi \in \mathbb{R}_+$. Then, define the set of $i=1, ..., m$ \textbf{delay-dependent} predictor operators mapping from $\mathbb{R}^n \times C^1([0,1];\mathbb{R})^{m-i+1} \times \mathbb{R}_+$ to $C^1([0,1];\mathbb{R}^n)$ taking
%  inputs $(Q,U_i,\dots,U_m, \varphi)$ and returning a function $P_i \in C^1([0,1];\mathbb{R}^n)$ where the output $P_i$ satisfies the implicit differential equation
% \begin{align}
%     P_i(s) - Q - \varphi \int_0^s f(P_i(\theta), \kappa_1(P_i(\theta)), \dots, \kappa_{i-1}(P_i(\theta)), U_i(\theta), \dots, U_m(\theta)) d\theta =& 0, \quad s \in [0,1], \label{eq:predictor_def_Pi}
% \end{align}
% where $\kappa_i$, $i \in [m-1]$, are a set of $C^1$ functions from $\mathbb{R}^n \to \mathbb{R}$ that are assumed to be apriori known as static parameters.
% \end{definition}
% Notice that, by definition, the predictor operator yields the solution to \eqref{eqs:ps} scaled to the interval $[0,1]$ where $\varphi_i$, $i \in [m]$, are delay values that are passed as inputs to the corresponding operators. That is,
% \begin{subequations}
% \begin{align}
%     P_1(t) =& \mathcal{P}_1(X(t), T_{D_1,0}(t)U_1, ..., T_{D_m,D_{m-1}}(t)U_m, D_1)(1), \\
%     P_2(t) =& \mathcal{P}_2(P_1(t), T_{D_{2,1},0}(t)U_2, ..., T_{D_{m,1},D_{m,2}}(t)U_m, D_{2,1})(1), \\
%     &\ \ \vdots \nonumber \\
%     P_m(t) =& \mathcal{P}_m(P_{m-1}(t), T_{D_{m,m-1},0}(t)U_m, D_{m,m-1})(1),
% \end{align}
% \end{subequations}
% for all $s \in [0,1]$, where $T_{\varphi_i,\varphi_j}$, $i, j \in [m]$, $i \geq j$, are the operators that, when applied to the control history $U$, yield
% \begin{align}
%     (T_{\varphi_i,\varphi_j}(t)U) := U(t-\varphi_i+(\varphi_i-\varphi_j)x), \quad x \in [0,1).
% \end{align}
% % In our investigation, we explored various definitions of predictor operators, including an alternative formulation with a single operator for the entire prediction horizon. However, we ultimately opted for the current approach, which utilizes $m$ separate operators. This choice was driven by the trade-off between differentiability of the operator on the delay boundaries and achieving a more concise and practical definition. Here, we ensure computational efficiency and scalability while still preserving the stability and performance guarantees of the predictor feedback system.
% The deliberate choice to use $m$ predictor operators instead of a single large operator offers several advantages. By doing so, the user can train $m$ distinct neural networks, each tailored to a specific dataset, rather than developing a monolithic model that must capture all $m$ prediction tasks from a single dataset. This modular approach simplifies training and allows inference to be parallelized across multiple compute nodes, leading to significant speedups compared to a single, unified operator. 
%  However, this design induces a trade-off: the lack of continuity at the boundaries between adjacent operators. Specifically, the outputs may not align perfectly at the interfaces, i.e., 
% $P_{j}(\cdot)(1) \neq P_{j+1}(\cdot)(0)$, 
% $j \in [m-1]$. Fortunately, this issue can be addressed in practice by enforcing continuity through network design by deliberately setting $P_{j+1}(\cdot)(0) = P_{j}(\cdot)(1)$ inside the neural operator architecture. 
% Now, for pedagogical reasons, we present the simplest possible two-input example illustrating how each neural operator is deployed, see Figure \ref{fig:block_diagram}.
% \begin{example}
% Consider a nonlinear system with two inputs under distinct input delays $0 < D_1 < D_2<\infty$
% \begin{align}
%     \dot{X}(t) = f(X(t), U_1(t-D_1), U_2(t-D_2))\,. 
% \end{align}
% The stabilizing control law is given by
% \begin{subequations}
% \begin{align}
%     U_1(t) =& \kappa_1(X(t+D_1)) = \kappa_1(P_1(t))\,, \label{eq:predictor_example_U1} \\ 
%     U_2(t) =& \kappa_2(X(t+D_2)) = \kappa_2(P_2(t)) \,, \label{eq:predictor_example_U2}
% \end{align}
% \end{subequations}
% where
% \begin{subequations}
% \begin{align}
%     P_1(t) =& X(t) + \int_{t-D_1}^t f(X(s), U_1(s), U_2(s-D_{2, 1})) ds\,, \label{eq:predictor_example_P1} \\
%     P_2(t) =& P_1(t) + \int_{t-{D_{2, 1}}}^t f(P_2(s), \kappa_1(P_2(s)), U_2(s)) ds\,. \label{eq:predictor_example_P2}
% \end{align}
% \end{subequations}
% In the operator reformulation, we have that $P_1, P_2$ are equivalent to
% \begin{subequations}
% \begin{align}
%     P_1(t) =& \mathcal{P}_1(X(t), T_{D_1,0}(t)U_1, T_{D_2,D_{2,1}}(t)U_2, D_1)(1), \\
%     P_2(t) =& \mathcal{P}_2(P_1(t), T_{D_{2,1},0}(t)U_2, D_{2,1})(1).
% \end{align}
% \end{subequations}
% \end{example}

% \begin{figure}[ht]
%     \centering
%     \includegraphics[width=\textwidth]{sections/section2/images/Multi-Input-Predictor-Feedback.png}
%     \caption{Block diagram of predictor-feedback system with neural approximated predictor operators with two distinct delays $0 < D_1 < D_2$.}
%     \label{fig:block_diagram}
% \end{figure}

% First, note that the predictors $P_1(t)$ and $P_2(t)$, although approximated for the entire horizon from $s \in [0, 1]$ are only needed in deployment for their final value. Additionally, we emphasize that, through the operators modular design, only the minimal amount of information as controlled by the deliberate shift $T_{\varphi_1, \varphi_2}(t)$ is needed as input to the operator hence limiting the amount of computational processing at each step. 

% We are now ready to discuss the existence of approximations to the operators $\mathcal{P}_i$ and highlight that neural operators, one particular class of neural network designs, exist as universal approximators of the predictor. 


% \subsection{Mathematical properties of the neural operator}
% In this section, our aim is to show that for any desired accuracy $\epsilon > 0$, there exist neural approximations of the predictor operator $\mathcal{P}_i$, $i \in [m]$.To do so, we aim to invoke the universal approximation operator theorem \cite{LLS2024}, and therefore, first prove the set of predictor operators are continuous. As a natural assumption for showing continuity of the predictor operator (See \cite{Luke_Peijia:2025}), we require local Lipschitz continuity of the dynamics: 

% \begin{assumption}\label{ass:Assumption4}
% Let $f(X,U_1,\dots,U_m)$ be as in \eqref{eq:system} and $\mathcal{X} \subset \mathbb{R}^n$, $\mathcal{U} \subset \mathbb{R}$ be two compact sets. Then, there exists a constant $C_f > 0$ such that $f$ satisfies the Lipschitz condition
% \begin{align}
%     \bigg| f(x_1,u_1,\dots,u_m)-f(x_2,u_{m+1},\dots,u_{2m}) \bigg| \leq C_f \bigg( |x_1-x_2| + \sum_{i=1}^m |u_i - u_{m+i}| \bigg),
% \end{align}
% for all $x_1,x_2 \in \mathcal{X}$ and $u_1,\dots,u_{2m} \in \mathcal{U}$.
% \end{assumption}
% Now, we prove the Lipschitz continuity of the predictor operator in the following lemma.

% \begin{lemma}
% Let $\mathcal{X}$ and $\mathcal{U}$ be compact sets and $C_f$ be the Lipschitz constant as in Assumption~\ref{ass:Assumption4}. Let $X_1, X_2 \in \mathcal{X}$, control functions $U_j, U_{m+j} \in C^1([0,1];\mathcal{U})$ with $j = i,\dots,m$, and positive real numbers $\varphi_1, \varphi_2 \in \Phi \subset \mathbb{R}_+$, such that $|X_1|,|X_2|<\overline{\mathcal{X}}$, $\|U_j\|_{L^\infty}, \|U_{m+j}\|_{L^\infty} < \overline{\mathcal{U}}$ with $j = i,\dots,m$, and $|\varphi_1|, |\varphi_2| < \overline{\Phi}$. Then, each predictor operator $\mathcal{P}_i$, $i \in [m]$, from Definition~\ref{def:delay_dependent_predictor_operators} satisfies the following Lipschitz condition
% \begin{align}
%     \bigg\| \mathcal{P}_i(X_1,U_i,\dots,U_m,\varphi_1) - \mathcal{P}_i(X_2,U_{m+i},\dots,U_{2m},\varphi_2) \bigg\|_{L^\infty} 
%     \leq& \nonumber \\
%     C_{\mathcal{P}} \Biggl( |X_1 - X_2| + \sum_{j=i}^m \bigg\| U_j - U_{m+j} \bigg\|_{L^\infty} 
%     &+ |\varphi_1 - \varphi_2| \Biggr), \label{eq:predictor_lipschitz_inequality}
% \end{align}
% with Lipschitz constant
% \begin{align}
%     C_{\mathcal{P}} := \max \bigg( 1, \Xi, \overline{\Phi} C_f \bigg) e^{\overline{\Phi} C_{\kappa}},
% \end{align}
% where $\Xi$ and $C_{\kappa}$ are constants defined by
% \begin{align}
%     \Xi &:= m C_f \overline{\mathcal{U}} + C_{\kappa} \bigg( \overline{\mathcal{X}} + m \overline{\Phi} C_f \overline{\mathcal{U}} \bigg) e^{\overline{\Phi} C_{\kappa}}, \\
%     C_{\kappa} &:= C_f \bigg( 1 + \sum_{i=1}^m C_i \bigg),
% \end{align}
% and where $C_i$ is Lipschitz constant of each control law $\kappa_i$, $i \in [m]$.
% \end{lemma}
% \begin{proof}
% Let define functions
% \begin{align}
%     P_i^{(1)}(s) &:= \mathcal{P}_i(X_1,U_i,\dots,U_m, \varphi_1)(s) \\
%     P_i^{(2)}(s) &:= \mathcal{P}_i(X_2,U_{m+i},\dots,U_{2m}, \varphi_2)(s),
% \end{align}
% for all $s \in [0,1]$. Before we proceed to prove the main part of the lemma, notice that each predictor is uniformly bounded in its domain. First, by definition and Assumption \ref{ass:Assumption4}, we have
% \begin{align}
%     \bigg| P_i^{(1)}(s) \bigg| \leq& |X_1| + \varphi_1 \int_{0}^s \bigg| f \bigg( P_i^{(1)}(\theta), \kappa_1 \bigg( P_i^{(1)}(\theta) \bigg), \dots, \kappa_{i-1} \bigg( P_i^{(1)}(\theta) \bigg), U_i(\theta), \dots, U_m(\theta) \bigg) \bigg| d\theta \nonumber \\
%         \leq& |X_1| + \varphi_1 \int_{0}^s C_f \Biggl( \bigg| P_i^{(1)}(\theta) \bigg| + \sum_{j=1}^{i-1} \bigg| \kappa_j \bigg( P_i^{(1)}(\theta) \bigg) \bigg| + \sum_{j=i}^m \bigg| U_j(\theta) \bigg| \Biggr) d\theta \nonumber \\
%         \leq& |X_1| + \varphi_1 \int_{0}^s C_f \Biggl( \bigg| P_i^{(1)}(\theta) \bigg| + \sum_{j=1}^{i-1} C_j \bigg| P_i^{(1)}(\theta) \bigg| + \sum_{j=i}^m \bigg| U_j(\theta) \bigg| \Biggr) d\theta \nonumber \\
%         \leq& |X_1| + \varphi_1 C_f \sum_{j=i}^m \| U_j \|_{L^\infty} + \varphi_1 C_{\kappa} \int_{0}^s \bigg| P_i^{(1)}(\theta) \bigg| d\theta,
% \end{align}
% where $C_{\kappa}$ is defined by
% \begin{align}
%     C_{\kappa} := C_f \bigg( 1 + \sum_{i=1}^m C_i \bigg).
% \end{align}
% Applying Gronwall's inequality yields
% \begin{align}
%     \bigg| P_i^{(1)}(s) \bigg| \leq& \bigg( \overline{\mathcal{X}} + m \overline{\Phi} C_f \overline{\mathcal{U}} \bigg) e^{\overline{\Phi} C_{\kappa}}.
% \end{align}
% Then, we have
% \begin{align}
%     \bigg| P_i^{(1)}(s) - P_i^{(2)}(s) \bigg| \leq& |X_1 - X_2| \nonumber \\
%             &\quad + \int_{0}^s \Bigl| \varphi_1 f \bigg( P_i^{(1)}(\theta), \kappa_1 \bigg( P_i^{(1)}(\theta) \bigg), \dots, \kappa_{i-1} \bigg( P_i^{(1)}(\theta) \bigg), U_i(\theta), \dots, U_m(\theta) \bigg) \nonumber \\
%             &\qquad - \varphi_2 f \bigg( P_i^{(2)}(\theta), \kappa_1 \bigg( P_i^{(2)}(\theta) \bigg), \dots, \kappa_{i-1} \bigg( P_i^{(2)}(\theta) \bigg), U_{m+i}(\theta), \dots, U_{2m}(\theta) \bigg) \Bigr| d\theta \nonumber \\
%         \leq& |X_1 - X_2| \nonumber \\
%             &\quad + |\varphi_1 - \varphi_2| \int_{0}^s \bigg| f \bigg( P_i^{(1)}(\theta), \kappa_1 \bigg( P_i^{(1)}(\theta) \bigg), \dots, \kappa_{i-1} \bigg( P_i^{(1)}(\theta) \bigg), U_i(\theta), \dots, U_m(\theta) \bigg) \bigg| d\theta \nonumber \\
%             &\quad + \varphi_2 \int_{0}^s \Bigl| f \bigg( P_i^{(1)}(\theta), \kappa_1 \bigg( P_i^{(1)}(\theta) \bigg), \dots, \kappa_{i-1} \bigg( P_i^{(1)}(\theta) \bigg), U_i(\theta), \dots, U_m(\theta) \bigg) \nonumber \\
%             &\qquad - f \bigg( P_i^{(2)}(\theta), \kappa_1 \bigg( P_i^{(2)}(\theta) \bigg), \dots, \kappa_{i-1} \bigg( P_i^{(2)}(\theta) \bigg), U_{m+i}(\theta), \dots, U_{2m}(\theta) \bigg) \Bigr| d\theta \nonumber \\
%         \leq& |X_1 - X_2| \nonumber \\
%             &\quad + |\varphi_1 - \varphi_2| C_f \sum_{j=i}^m \| U_j \|_{L^\infty} + |\varphi_1 - \varphi_2| C_{\kappa} \int_{0}^s \bigg| P_i^{(1)}(\theta) \bigg| d\theta \nonumber \\
%             &\quad + \varphi_2 \int_{0}^s C_f \Biggl( \bigg| P_i^{(1)}(\theta) - P_i^{(2)}(\theta) \bigg| \nonumber \\
%             &\qquad  + \sum_{j=1}^{i-1} \bigg| \kappa_j \bigg( P_i^{(1)}(\theta) \bigg) - \kappa_j \bigg( P_i^{(2)}(\theta) \bigg) \bigg| + \sum_{j=i}^m \bigg| U_j(\theta) - U_{m+j}(\theta) \bigg| \Biggr) d\theta \nonumber \\
%         \leq& |X_1 - X_2| \nonumber \\
%             &\quad + |\varphi_1 - \varphi_2| \bigg[ m C_f \overline{\mathcal{U}} + C_{\kappa} \bigg( \overline{\mathcal{X}} + m \overline{\Phi} C_f \overline{\mathcal{U}} \bigg) e^{\overline{\Phi} C_{\kappa}} \bigg] \nonumber \\
%             &\quad + \varphi_2 \int_{0}^s C_f \Biggl( \bigg| P_i^{(1)}(\theta) - P_i^{(2)}(\theta) \bigg| \nonumber \\
%             &\qquad  + \sum_{j=1}^{i-1} C_j \bigg( P_i^{(1)}(\theta) - P_i^{(2)}(\theta) \bigg) + \sum_{j=i}^m \bigg| U_j(\theta) - U_{m+j}(\theta) \bigg| \Biggr) d\theta \nonumber \\
%         \leq& |X_1 - X_2| \nonumber \\
%             &\quad + |\varphi_1 - \varphi_2| \bigg[ m C_f \overline{\mathcal{U}} + C_{\kappa} \bigg( \overline{\mathcal{X}} + m \overline{\Phi} C_f \overline{\mathcal{U}} \bigg) e^{\overline{\Phi} C_{\kappa}} \bigg] \nonumber \\
%             &\quad + \varphi_2 C_f \sum_{j=1}^m \| U_j - U_{m+j} \|_{L^\infty} \nonumber \\
%             &\qquad + \varphi_2 C_{\kappa} \int_0^s \bigg| P_i^{(1)}(\theta) - P_i^{(2)}(\theta) \bigg| d\theta, \label{eq:setup_for_gronwall}
% \end{align}
% Applying Gronwall's inequality again yields the desired result.
% \end{proof}

% Applying the continuity of $\mathcal{P}_i$, $i \in [m]$, with the universal approximation theorem of neural operators yields
% \begin{theorem}\label{thm:predictor_operator_theorem}
% Let $Q \in \mathcal{X}$, $U_i \in C^1([0,1];\mathcal{U})$, $i \in [m]$, and $\varphi \in \Phi$, where $\mathcal{X} \subset \mathbb{R}^n$, $\mathcal{U} \subset \mathbb{R}$, and $\Phi \subset \mathbb{R}_+$ are bounded domains. For each $i \in [m]$, fix a compact set $K_i \subset \mathcal{X} \times C^1([0,1];\mathcal{U})^{m-i+1} \times \Phi$. Then, for any $\epsilon > 0$, there exist neural operator approximations $\hat{\mathcal{P}}_i: K_i \to C^1([0,1];\mathbb{R}^n)$ of the predictor operators $\mathcal{P}_i$ from Definition \ref{def:delay_dependent_predictor_operators} such that
% \begin{equation}
% \sup_{(Q,U_i,\ldots,U_m,\varphi) \in K_i} \|\mathcal{P}_i(Q,U_i,\ldots,U_m,\varphi)(s) - \hat{\mathcal{P}}_i(Q,U_i,\ldots,U_m,\varphi)(s)\|_{L^\infty} < \epsilon, \quad \forall s \in [0,1],
% \end{equation}
% for all $X \in \mathcal{X}$, $U_i \in C^1([0,1);\mathcal{U})$, $i \in [m]$, and $\varphi \in \Phi$.
% \end{theorem}

% It is important to notice that the theorem restricts the possibility of global stability guarantees due to the compactness of intervals on which it holds. A similar type of restriction appears in \cite{Luke_Peijia:2025} and is a common restriction of almost all universal approximation theorems. This restriction directly results in the semiglobal stability estimate, as the initial conditions must be constrained in order for the neural operator approximation to hold for all $t \geq 0$. Theorem \ref{thm:predictor_operator_theorem} guarantees the existence of neural approximated predictor operators, but does not provide a minimum number of network parameters or data samples to achieve such errors.

% \subsection{Perturbed Transport PDE under the approximated predictor}
% We now focus on analyzing the resulting stability properties of the system \eqref{eq:system} with the approximate predictors. For notational simplicity, define $\hat{P}_i(t)$, $i \in [m]$, as the solutions to the neural operator approximations for all $t \geq 0$, namely
% \begin{subequations}
% \begin{align}
%     \hat{P}_1(t) =& \hat{\mathcal{P}}_1(X(t), T_{D_1,0}(t)U_1, ..., T_{D_m,D_{m-1}}(t)U_m, D_1)(1), \\
%     \hat{P}_2(t) =& \hat{\mathcal{P}}_2(\hat{P}_1(t), T_{D_{2,1},0}(t)U_2, ..., T_{D_{m,1},D_{m,2}}(t)U_m, D_{2,1})(1), \\
%     &\ \ \vdots \nonumber \\
%     \hat{P}_m(t) =& \hat{\mathcal{P}}_m(\hat{P}_{m-1}(t), T_{D_{m,m-1},0}(t)U_m, D_{m,m-1})(1).
% \end{align}
% \end{subequations}
% Instead of employing exact predictors in predictor-feedback systems, we propose using neural operator-approximated predictors. Thus, the plant can be represented as
% \begin{subequations}\label{eqs:plant}
%     \begin{align}
%         \dot{X}(t) =& f \bigg( X(t), u_1(0,t), \dots, u_m(0,t) \bigg), \quad t \in [0,\infty) \label{eq:plant_ode} \\
%         \partial_t u_i(x,t) =& \partial_x u_i(x,t), \quad (x,t) \in (0,D_i) \times [0,\infty) \label{eq:plant_pde} \\
%         u_i(D_i,t) =& \kappa_i(\hat{P}_i(t)), \quad t \in [0,\infty), \label{eq:plant_pde_boundary}
%     \end{align}
% \end{subequations}
% for all $i \in [m]$. The following part of this section includes a series of technical lemmas presented next, which are required to prove the main result of this work. Lastly, we assume the existence of globally asymptotically stabilizing control laws for the system without delays. The assumption can be relaxed with the expectation that the stability result for the predictor feedback controller correspondingly weakens.



% \begin{lemma}\label{lem:lemma_1_of_6}
% The backstepping transformations of $u_i$, for all $i \in [m]$, defined by
% \begin{align}
%     w_1(x,t) =& u_1(x,t) - \kappa_1\bigg(p_1(x,t)\bigg), \quad x \in [0, D_1] \label{eq:w1} \\
%     w_2(x,t) =& u_2(x,t) - 
%         \begin{cases} 
%             \kappa_2\bigg(p_1(x,t)\bigg), \quad x \in [0, D_1] \\
%             \kappa_2\bigg(p_2(x,t)\bigg), \quad x \in [D_1, D_2]
%         \end{cases} \label{eq:w2} \\
%     &\ \vdots \nonumber \\
%     w_m(x,t) =& u_m(x,t) - 
%         \begin{cases} 
%             \kappa_m\bigg(p_1(x,t)\bigg), \quad x \in [0, D_1] \\
%             \kappa_m\bigg(p_2(x,t)\bigg), \quad x \in [D_1, D_2] \\
%             \ \vdots \\
%             \kappa_m\bigg(p_m(x,t)\bigg), \quad x \in [D_{m-1}, D_m]
%         \end{cases}, \label{eq:wm}
% \end{align}
% where $p_i$, $i \in [m]$, are defined in \eqref{eqs:ps}, together with the control laws \eqref{eq:control_laws}, transform the plant \eqref{eqs:plant} to the target system
% \begin{subequations}\label{eqs:target}
%     \begin{align}
%         \dot{X}(t) =& f \bigg( X(t), \kappa_1(X(t)) + w_1(0,t), \dots, \kappa_m(X(t)) + w_m(0,t) \bigg), \quad t \in [0,\infty) \label{eq:target_ode} \\
%         \partial_t w_i(x,t) =& \partial_x w_i(x,t), \quad (x,t) \in (0,D_i) \times [0,\infty) \label{eq:target_pde} \\
%         w_i(D_i,t) =& \kappa_i(\hat{P}_i(t)) - \kappa_i(P_i(t)), \quad t \in [0,\infty), \label{eq:target_pde_boundary}
%     \end{align}
% \end{subequations}
% for all $i \in [m]$.
% \end{lemma}
% \begin{proof}
% Setting $x = 0$ into \eqref{eq:w1}-\eqref{eq:wm} and using \eqref{eq:p1(0,t)}, we obtain \eqref{eq:target_ode}. Next, if we take both partial derivatives, $\partial_t$ and $\partial_x$, of \eqref{eq:w1}-\eqref{eq:wm}, we get
% \begin{align}
%     \partial_tw_1(x,t) =& \partial_tu_1(x,t) - \kappa_1'\bigg(p_1(x,t)\bigg)\partial_t\bigg(p_1(x,t)\bigg), \quad x \in (0, D_1) \\
%     \partial_xw_1(x,t) =& \partial_xu_1(x,t) - \kappa_1'\bigg(p_1(x,t)\bigg)\partial_x\bigg(p_1(x,t)\bigg), \quad x \in (0, D_1),
% \end{align}
% and
% \begin{align}
%     \partial_tw_i(x,t) =& \partial_tu_i(x,t) - 
%         \begin{cases} 
%             \kappa_i'\bigg(p_1(x,t)\bigg)\partial_t\bigg(p_1(x,t)\bigg), \quad x \in (0, D_1) \\
%             \kappa_i'\bigg(p_2(x,t)\bigg)\partial_t\bigg(p_2(x,t)\bigg), \quad x \in (D_1, D_2) \\
%             \ \vdots \\
%             \kappa_i'\bigg(p_i(x,t)\bigg)\partial_t\bigg(p_i(x,t)\bigg), \quad x \in (D_{i-1}, D_i)
%         \end{cases} \\
%     \partial_xw_i(x,t) =& \partial_xu_i(x,t) - 
%         \begin{cases} 
%             \kappa_i'\bigg(p_1(x,t)\bigg)\partial_x\bigg(p_1(x,t)\bigg), \quad x \in (0, D_1) \\
%             \kappa_i'\bigg(p_2(x,t)\bigg)\partial_x\bigg(p_2(x,t)\bigg), \quad x \in (D_1, D_2) \\
%             \ \vdots \\
%             \kappa_i'\bigg(p_i(x,t)\bigg)\partial_x\bigg(p_i(x,t)\bigg), \quad x \in (D_{i-1}, D_i)
%         \end{cases},
% \end{align}
% $i = 2, \dots, m$. In addition, partial derivatives exist on the boundaries of domains due to the continuous differentiability of functions $w_i(x,t)$, i.e. $\lim_{x \to D_j^-} \partial_tw_i(x,t) = \lim_{x \to D_j^+} \partial_tw_i(x,t)$ and $\lim_{x \to D_j^-} \partial_xw_i(x,t) = \lim_{x \to D_j^+} \partial_xw_i(x,t)$ for all $i = 2, \dots, m$, $j = 1, \dots, i-1$. Noting that $\partial_tu_i(x,t) = \partial_xu_i(x,t)$, holds for all $i \in [m]$, \eqref{eq:plant_pde}, and that the functions $p_i(x,t)$, $i \in [m]$, are actually functions of one variable, meaning that $\partial_tp_i(x,t) = \partial_xp_i(x,t)$, for all $x \in (0, D_i)$, $t > 0$, we obtain \eqref{eq:target_pde}. Lastly, we obtain \eqref{eq:target_pde_boundary} by setting $x = D_i$ into \eqref{eq:w1}-\eqref{eq:wm}.
% \end{proof}


% \begin{lemma}\label{lem:lemma_2_of_6}
% The inverse backstepping transformations of \eqref{eq:w1}-\eqref{eq:wm} are defined by
% \begin{align}
%     u_1(x,t) =& w_1(x,t) + \kappa_1\bigg(\pi_1(x,t)\bigg), \quad x \in [0, D_1] \label{eq:inv_w1} \\
%     u_2(x,t) =& w_2(x,t) +
%         \begin{cases}
%             \kappa_2\bigg(\pi_1(x,t)\bigg), \quad x \in [0, D_1] \\
%             \kappa_2\bigg(\pi_2(x,t)\bigg), \quad x \in [D_1, D_2]
%         \end{cases} \label{eq:inv_w2} \\
%     &\ \vdots \nonumber \\
%     u_m(x,t) =& w_m(x,t) -
%         \begin{cases}
%             \kappa_m\bigg(\pi_1(x,t)\bigg), \quad x \in [0, D_1] \\
%             \kappa_m\bigg(\pi_2(x,t)\bigg), \quad x \in [D_1, D_2] \\
%             \ \vdots \\
%             \kappa_m\bigg(\pi_m(x,t)\bigg), \quad x \in [D_{m-1}, D_m]
%         \end{cases}, \label{eq:inv_wm}
% \end{align}
% where
% \begin{align}
%     \pi_1(x,t) =& X(t) + \int_0^x f \bigg (\pi_1(y,t), w_1(y,t) + \kappa_1 \bigg( \pi_1(y,t) \bigg), \ldots, w_m(y,t) + \kappa_m \bigg( \pi_1(y,t) \bigg) \bigg) dy, \label{eq:pi1} \\
%         &\qquad x \in [0, D_1] \nonumber \\
%     \pi_i(x,t) =& \pi_{i-1}(D_{i-1},t) + \int_{D_{i-1}}^x f \bigg( \pi_i(y,t), w_1(y,t) + \kappa_1 \bigg( \pi_i(y,t) \bigg), \ldots, w_m(y,t) + \kappa_m \bigg( \pi_m(y,t) \bigg) \bigg) dy, \label{eq:pii} \\
%         &\qquad x \in [D_{i-1}, D_i], \quad i=2,\dots,m. \nonumber
% \end{align}
% \end{lemma}
% \begin{proof}
% We prove this lemma by showing that $p_1(x,t) = \pi_1(x,t)$, for all $x \in [0,D_1]$ and $p_i(x,t) = \pi_i(x,t)$, for all $x \in [D_{i-1},D_i]$, $i=2,\dots,m$. We observe that \eqref{eq:pi1} can be equivalently written as an initial value problem
% \begin{align}
%     \partial_x \pi_1(x,t) =& f \bigg( \pi_1(x,t), w_1(x,t) + \kappa_1(\pi_1(x,t)), \ldots, w_m(x,t) + \kappa_m(\pi_1(x,t)) \bigg) \label{eq:pi1_ode_x} \\
%     \pi_1(0,t) =& X(t) \label{eq:pi1_ode_init}
% \end{align}
% for all $x \in (0,D_1)$, whose solution is given as $\pi_1(x,t) = X(t+x)$. Firstly, the solution satisfies the boundary condition \eqref{eq:pi1_ode_init}. Secondly, it also satisfies the ODE \eqref{eq:pi1_ode_x} where by substituting $\pi_1(x,t) = X(t+x)$, we get
% \begin{align}
%     X'(t+x) = f \bigg( X(t+x), w_1(0,t+x) + \kappa_1(X(t+x)), \dots, w_m(0,t+x) + \kappa_m(X(t+x)) \bigg), \label{eq:lemma2_X1_prime}
% \end{align}
% for all $x \in (0,D_1)$ and $t \geq 0$, where we used the fact that $w_i(x,t)$, $i \in [m]$, are functions of one variable, namely $w_i(x,t) = w_i(0,t+x)$. Applying the substitution $s=t+x$ in \eqref{eq:lemma2_X1_prime}, we obtain \eqref{eq:target_ode}, and therefore we prove that $X(t+x)$ is a solution to \eqref{eq:pi1_ode_x}, \eqref{eq:pi1_ode_init}, for all $x \in [0,D_1]$. In addition, since $f$ is a locally Lipschitz vector field, $\pi_1(x,t) = X(t+x)$ is the only solution. Using \eqref{eq:P_ip_i}, we arrive at $p_1(x,t)=\pi_1(x,t)$, for all $x \in [0,D_1]$.

% Assume now that $\pi_j(x,t)=X(t+x)$, for all $x \in [D_{j-1},D_j]$, and some $j$. Then, by induction, we claim that $\pi_{j+1}=X(t+x)$, for all $x \in [D_j,D_{j+1}]$. The function $\pi_{j+1}(x,t)$, based on \eqref{eq:pii}, satisfies the following initial value problem
% \begin{align}
%     \partial_x \pi_{j+1}(x,t) =& f \bigg( \pi_{j+1}(x,t), w_1(x,t) + \kappa_1(\pi_{j+1}(x,t)), \ldots, w_m(x,t) + \kappa_m(\pi_{j+1}(x,t)) \bigg) \label{eq:pij+1_ode_x} \\
%     \pi_{j+1}(D_j,t) =& \pi_j(D_j,t), \label{eq:pij+1_ode_init}
% \end{align}
% for all $x \in [D_j,D_{j+1}]$. Analogously to the base case, we conclude that the solution $\pi_{j+1}(x,t)=X(t+x)$, for all $x \in [D_j,D_{j+1}]$, satisfies both the boundary condition \eqref{eq:pij+1_ode_init} and the ODE \eqref{eq:pij+1_ode_x}, where by substitution, we get
% \begin{align}
%     X'(t+x) =& f \bigg( X(t+x), w_1(0,t+x) + \kappa_1(X(t+x)), \ldots, w_m(0,t+x) + \kappa_m(X(t+x)) \bigg),
% \end{align}
% for all $x \in [D_j,D_{j+1}]$, $t \geq 0$. Again, we used the fact that $w_i(x,t)$, $i \in [m]$, are functions of one variable, namely $w_i(x,t) = w_i(0,t+x)$. Analogously, we arrive at the same conclusion that $p_j(x,t)=\pi_j(x,t)$, for all $x \in [D_{j-1},D_j]$. Since this holds for an arbitrary $j$, the proof is complete.
% \end{proof}


% \begin{lemma}\label{lem:lemma_3_of_6}
% Let the system \eqref{eqs:target} satisfy the Assumption \ref{ass:Assumption2}. Then, there exists a class $\mathcal{KL}$ function $\beta_1$ and a class $\mathcal{K}$ function $\gamma_1$ such that the following holds
% \begin{align}
%     \breve{\Gamma}(t) \leq& \beta_1 \bigg( \breve{\Gamma}(0), t \bigg) + \gamma_1 \bigg( \sup_{0 \leq \tau \leq t} \bigg\{ \sum_{i=1}^m |\kappa_i(\hat{P}(-D_{m,i},t)) - \kappa_i(P_i(t))| \bigg\} \bigg), \quad t \geq 0, \label{eq:lemma3_Gamma_bar}
% \end{align}
% where
% \begin{align}
%     \breve{\Gamma}(t) =& |X(t)| + \sum_{i=1}^m \|w_i(t)\|_{L^\infty}. \label{eq:Gamma_bar_def}
% \end{align}
% \end{lemma}
% \begin{proof}
% From the Assumption \ref{ass:Assumption2} and the definition of input-to-state stability in \cite{SontagISS:1995}, it follows that there exist function $\beta_2 \in \mathcal{KL}$ and $\gamma_2 \in \mathcal{K}$ such that
% \begin{align}
%     |X(t)| \leq& \beta_2(|X(0)|, t) + \gamma_2 \bigg( \sup_{0 \leq \tau \leq t} \bigg\{ \sum_{i=1}^m |w_i(0,\tau)| \bigg\} \bigg) \label{eq:w_iss}
% \end{align}
% for all $t \geq 0$. Furthermore, we have the ISS estimate for $\mathcal{L}^\infty$ of perturbed transport PDE from \cite{Luke_Peijia:2025} (Lemma 6) for some constant $c > 0$,
% \begin{align}
%     \|w_i(t)\|_{L^\infty} \leq& \exp(c(D_i-t)) \|w_i(0)\|_{L^\infty} + \exp(cD_i) \sup_{0 \leq \tau \leq t} |w_i(D_i,\tau)|, \label{eq:wi_pde_iss}
% \end{align}
% for all $i \in [m]$. The following calculations are the setup for Lemma C.4 in \cite{KKK}. From \eqref{eq:w_iss}, we have
% \begin{align}
%     |X(t)| \leq& \beta_2(|X(0)|, t) + \gamma_2 \bigg( \sup_{0 \leq \tau \leq t} \bigg\{ \sum_{i=1}^m |w_i(0,\tau)| \bigg\} \bigg) \nonumber \\
%         \leq& \beta_2(|X(0)|, t) + \gamma_2 \bigg( \sup_{0 \leq \tau \leq t} \bigg\{ \sum_{i=1}^m \|w_i(\tau)\|_{L^\infty} \bigg\} \bigg) \nonumber \\
%         \leq& \beta_2(|X(0)|, t) + \gamma_2 \bigg( \sup_{0 \leq \tau \leq t} \bigg\{ \sum_{i=1}^m \|w_i(\tau)\|_{L^\infty} + \sum_{i=1}^m |w_i(D_i,\tau)| \bigg\} \bigg) \label{eq:lemma3_x_iss},
% \end{align}
% and, from \eqref{eq:wi_pde_iss} applied to all $i \in [m]$, we have
% \begin{align}
%     \sum_{i=1}^m \|w_i(t)\|_{L^\infty} \leq& \sum_{i=1}^m \exp(c(D_i-t)) \|w_i(0)\|_{L^\infty} + \sum_{i=1}^m \exp(cD_i) \sup_{0 \leq \tau \leq t} \bigg|w_i(D_i,\tau)\bigg| \nonumber \\
%         \leq& \beta_3 \bigg( \sum_{i=1}^m \|w_i(0)\|_{L^\infty}, t \bigg) + \gamma_3 \bigg( \sup_{0 \leq \tau \leq t} \bigg\{ \sum_{i=1}^m |w_i(D_i,\tau)| \bigg\} \bigg), \label{eq:lemma3_w_iss}
% \end{align}
% where $\beta_3 \in \mathcal{KL}$, and $\gamma_3 \in \mathcal{K}$. Now, by reproducing the argument from Lemma C.4 in \cite{KKK}, we obtain
% \begin{align}
%     \bigg| \bigg(X(t), \sum_{i=1}^m \|w_i(t)\|_{L^\infty} \bigg) \bigg| \leq& \beta_4 \bigg( \bigg| \bigg( X(0), \sum_{i=1}^m \|w_i(0)\|_{L^\infty} \bigg) \bigg|, t \bigg) \nonumber \\
%         &\qquad + \gamma_4 \bigg( \sup_{0 \leq \tau \leq t} \bigg\{ \sum_{i=1}^m |w_i(D_i,\tau)| \bigg\} \bigg), \label{eq:after_c4}
% \end{align}
% where $\beta_4 \in \mathcal{KL}$, and $\gamma_4 \in \mathcal{K}$. Furthermore, using the facts that $|x_1| \leq |(x_1,x_2)|$, $|x_2| \leq |(x_1,x_2)|$, and $|(x_1,x_2)| \leq |x_1| + |x_2|$, we obtain the following two inequalities
% \begin{align}
%     |X(t)| \leq& \beta_4 \bigg( |X(0)| + \sum_{i=1}^m \|w_i(0)\|_{L^\infty}, t \bigg) + \gamma_4 \bigg( \sup_{0 \leq \tau \leq t} \bigg\{ \sum_{i=1}^m |w_i(D_i,\tau)| \bigg\} \bigg), \label{eq:x_after_c4}
% \end{align}
% and
% \begin{align}
%     \sum_{i=1}^m \|w_i(0)\|_{L^\infty} \leq& \beta_4 \bigg( |X(0)| + \sum_{i=1}^m \|w_i(0)\|_{L^\infty}, t \bigg) + \gamma_4 \bigg( \sup_{0 \leq \tau \leq t} \bigg\{ \sum_{i=1}^m |w_i(D_i,\tau)| \bigg\} \bigg). \label{eq:w_after_c4}
% \end{align}
% It follows from \eqref{eq:x_after_c4} and \eqref{eq:w_after_c4} that
% \begin{align}
%     |X(t)| + \sum_{i=1}^m \|w_i(0)\|_{L^\infty} \leq& \beta_1 \bigg( |X(0)| + \sum_{i=1}^m \|w_i(0)\|_{L^\infty}, t \bigg) + \gamma_1 \bigg( \sup_{0 \leq \tau \leq t} \bigg\{ \sum_{i=1}^m |w_i(D_i,\tau)| \bigg\} \bigg), \label{eq:x_and_w_after_c4}
% \end{align}
% where $\beta_1 \in \mathcal{KL}$, and $\gamma_1 \in \mathcal{K}$. Lastly, by defining the function $\breve{\Gamma}(t) = |X(t)| + \sum_{i=1}^m \|w_i(t)\|_{L^\infty}$ and using \eqref{eq:target_pde_boundary}, we obtain the following bound
% \begin{align}
%     \breve{\Gamma}(t) \leq& \beta_1 \bigg( \breve{\Gamma}(0), t \bigg) + \gamma_1 \bigg( \sup_{0 \leq \tau \leq t} \bigg\{ \sum_{i=1}^m |\kappa_i(\hat{P}(-D_{m,i},t)) - \kappa_i(P_i(t))| \bigg\} \bigg),
% \end{align}
% for all $t \geq 0$, which completes the proof.
% \end{proof}

% \textcolor{red}{For example, here you give no guidance to the reader about why the proofs of lemma 13 and 14 are useful (I know you need them to transform back to original system, but the reader does not). Anyther thing is, are these lemmas already given in nikos paper - no need to reprove them here }

% \begin{lemma}\label{lem:lemma_4_of_6}
% There exist class $\mathcal{K}_{\infty}$ functions $\rho_{1},\ldots,\rho_{m}$ such that
% \begin{align}
%     \|p_{1}(t)\|_{\infty} \leq& \rho_{1}\bigg(\Xi(t)\bigg) \\
%     \|p_{2}(t)\|_{\infty} \leq& \rho_{2}\bigg(\Xi(t)\bigg) \\
%     &\ \vdots \nonumber \\
%     \|p_{m}(t)\|_{\infty} \leq& \rho_{m}\bigg(\Xi(t)\bigg),
% \end{align}
% where $\Xi$ is defined in \eqref{eq:Xi}.
% \end{lemma}


% \begin{lemma}\label{lem:lemma_5_of_6}
% Let the system \eqref{eqs:target} satisfy the Assumption \ref{ass:Assumption2}. There exist class $\mathcal{K}_\infty$ functions $\breve{\rho}_1, \dots, \breve{\rho}_m, $ such that
% \begin{align}
%     \|\pi_1(t)\|_\infty \leq& \breve{\rho}_1\bigg(\breve{\Gamma}(t)\bigg) \label{eq:pi1_bound} \\
%     \|\pi_2(t)\|_\infty \leq& \breve{\rho}_2\bigg(\breve{\Gamma}(t)\bigg) \label{eq:pi2_bound} \\
%     &\ \vdots \nonumber \\
%     \|\pi_m(t)\|_\infty \leq& \breve{\rho}_m\bigg(\breve{\Gamma}(t)\bigg), \label{eq:pim_bound}
% \end{align}
% where $\breve{\Gamma}(t)$ is defined in \eqref{eq:Gamma_bar_def}.
% \end{lemma}
% \begin{proof}
% We prove the lemma by induction. For clarity, we present the initial step and the induction step. First, we prove the bound \eqref{eq:pi1_bound}. Based on the Assumption \ref{ass:Assumption2}, the fact that the $\pi_1(x,t)$ satisfy the initial value problem \eqref{eq:pi1_ode_x}-\eqref{eq:pi1_ode_init}, and the definition of input-to-state stability in \cite{SontagISS:1995}, we have
% \begin{align}
%     |\pi_1(x,t)| \leq& \beta_5\bigg(|X(t)|, x\bigg) + \gamma_5 \bigg( \sup_{0 \leq y \leq x} \bigg( \sum_{i=1}^m |w_i(y,t)| \bigg) \bigg) \nonumber \\
%         \leq& \beta_5\bigg(|X(t)|, x\bigg) + \gamma_5 \bigg( \sum_{i=1}^m \sup_{0 \leq y \leq x} |w_i(y,t)| \bigg) \nonumber \\
%         \leq& \beta_5\bigg(|X(t)|, x\bigg) + \gamma_5\bigg(\sum_{i=1}^m \|w_i(t)\|_{L^\infty} \bigg) \nonumber \\
%         \leq& \beta_5\bigg(\breve{\Gamma}(t), x\bigg) + \gamma_5\bigg(\breve{\Gamma}(t)\bigg) \label{eq:lemma5_pi1_bound}
% \end{align}
% for all $x \in [0, D_1]$ and $t \geq 0$, where $\beta_5 \in \mathcal{KL}$, and $\gamma_5 \in \mathcal{K}$. Furthermore, we define function $\breve{\rho}_1 \in \mathcal{K}$, $\breve{\rho}_1(s) = \beta_5(s,0) + \gamma_5(s)$, and, by combining it with \eqref{eq:lemma5_pi1_bound}, we obtain the desired bound \eqref{eq:pi1_bound}. Now, assume that for some $j$ it holds that
% \begin{align}
%     \|\pi_j(t)\|_{L^\infty[0,D_j]} \leq \breve{\rho}_j(\breve{\Gamma}(t)), \quad t \geq 0.
% \end{align}
% Again, based on the Assumption \ref{ass:Assumption2}, the fact that the $\pi_{j+1}(x,t)$ satisfy the initial value problem \eqref{eq:pij+1_ode_x}-\eqref{eq:pij+1_ode_init}, and the definition of ISS, we have
% \begin{align}
%     |\pi_{j+1}(x,t)| \leq& \beta_6\bigg(|\pi_j(D_j)|, x-D_j\bigg) + \gamma_6\bigg(\sup_{D_j \leq y \leq x} \bigg(\sum_{i=1}^m |w_i(y,t)|\bigg)\bigg) \nonumber \\
%         \leq& \beta_6\bigg(|\pi_j(D_j)|, x-D_j\bigg) + \gamma_6\bigg(\sum_{i=1}^m \sup_{D_j \leq y \leq x} |w_i(y,t)| \bigg) \nonumber \\
%         \leq& \beta_6\bigg(|\pi_j(D_j)|, x-D_j\bigg) + \gamma_6\bigg(\sum_{i=1}^m \|w_i(t)\|_{L^\infty} \bigg) \nonumber \\
%         \leq& \beta_6\bigg(\breve{\rho}_j(\breve{\Gamma}(t)), x-D_j\bigg) + \gamma_6 \bigg( \breve{\Gamma}(t) \bigg) \label{eq:lemma5_pij+1_bound}
% \end{align}
% for all $x \in [D_j, D_{j+1}]$ and $t \geq 0$, where $\beta_6 \in \mathcal{KL}$, and $\gamma_6 \in \mathcal{K}$. Furthermore, we define function $\breve{\rho}_{j+1} \in \mathcal{K}$, $\breve{\rho}_{j+1}(s) = \beta_6(\breve{\rho}_j(s),0) + \gamma_6(s)$, and, by combining it with \eqref{eq:lemma5_pij+1_bound}, we obtain the following bound
% \begin{align}
%     \|\pi_{j+1}(t)\|_{L^\infty[0,D_j]} \leq \breve{\rho}_{j+1}(\breve{\Gamma}(t)),
% \end{align}
% for all $t \geq 0$.
% \end{proof}

% \begin{lemma}\label{lem:lemma_6_of_6}
% There exist class $\mathcal{K}_{\infty}$ functions $\rho$, $\breve{\rho}$ such that
% \begin{align}
% \breve{\Gamma}(t) \leq& \rho\bigg(\Gamma(t)\bigg) \label{eq:lemma6_Gamma_bar} \\
% \Gamma(t) \leq& \breve{\rho}\bigg(\breve{\Gamma}(t)\bigg), \label{eq:lemma6_Gamma}
% \end{align}
% where $\breve{\Gamma}$ is defined in \eqref{eq:Gamma_bar_def} and 
% \begin{align}
%     \Gamma(t) =& |X(t)| + \sum_{i=1}^m \|u_i(t)\|_{L^\infty}. \label{eq:Gamma_def}
% \end{align}
% \end{lemma}

% For conciseness, Lemma \eqref{lem:lemma_4_of_6} and Lemma \eqref{lem:lemma_6_of_6} are restated from \cite[Lemmas 4 and 6]{Bekiaris-Liberis_Krstic:2016} because the proof of the following main theorem relies on those two lemmas.

% \begin{theorem}\label{thm:main_result}
% Let the system \eqref{eqs:plant} satisfy Assumptions \ref{ass:Assumption1}-\ref{ass:Assumption3}, \ref{ass:Assumption4}, and assume that $\kappa_i$, $i \in [m]$, are stabilizing as in the Assumption \ref{ass:Assumption5}. Then, for $\overline{\mathcal{B}} := \min \{ \overline{\mathcal{X}}, \overline{\mathcal{U}} \}$, there exist functions $\alpha_1 \in \mathcal{K}$ and $\beta \in \mathcal{KL}$ such that if $\epsilon < \epsilon^* /m$ where
% \begin{align}
%     \epsilon^{*}(\overline{\mathcal{B}}) := \alpha_1^{-1}(\overline{\mathcal{B}}), \label{eq:epsilon_star}
% \end{align}
% and the initial state is constrained to
% \begin{align}
%     \Gamma(0) \leq \Omega, \label{eq:init_state_restriction}
% \end{align}
% where
% \begin{align}
%     \Omega(\epsilon, \overline{\mathcal{B}}) := \alpha_{2}^{-1} \bigg( \overline{\mathcal{B}} - \alpha_{1}(m\epsilon)\bigg), \quad \alpha_{2}(r) = \beta(r, 0), \ r \geq 0, \label{eq:Omega}
% \end{align}
% then, the closed-loop solutions, under the controller with the neural operator-approximated predictor $U_i(t) = \kappa_i \bigg( \hat{P}_i(t) \bigg)$, $i \in [m]$, satisfy the $\overline{\mathcal{B}}$-semiglobal and $\epsilon$-practical stability estimate
% \begin{align}
%     \Gamma(t) \leq& \beta(\Gamma(0), t) + \alpha_1(m\epsilon), \quad \forall t \geq 0, \label{eq:main_inequality}
% \end{align}
% where
% \begin{align}
%     \Gamma(t) =& |X(t)| + \sum_{i=1}^m \|u_i(t)\|_{L^\infty}, \label{eq:Xi}
% \end{align}
% with the neural operator semiglobally trained with $\overline{\mathcal{B}} > \alpha_2(\Omega)$ and $\epsilon \in (0, \alpha_1^{-1}(\overline{\mathcal{B}})/m)$.
% \end{theorem}
% \begin{proof}
% Combining \eqref{eq:lemma6_Gamma} with \eqref{eq:lemma3_Gamma_bar}, we get
% \begin{align}
%     \Gamma(t) \leq& \breve{\rho} \bigg( \beta_1 \bigg( \breve{\Gamma}(0), t \bigg) + \gamma_1 \bigg( \sup_{0 \leq \tau \leq t} \bigg\{ \sum_{i=1}^m |\kappa_i(\hat{P}_i(t)) - \kappa_i(P_i(t))| \bigg\} \bigg) \bigg) \nonumber \\
%         \leq& \breve{\rho} \bigg( 2\beta_1 \bigg( \breve{\Gamma}(0), t \bigg) \bigg) + \breve{\rho} \bigg( 2\gamma_1 \bigg( \sup_{0 \leq \tau \leq t} \bigg\{ \sum_{i=1}^m |\kappa_i(\hat{P}_i(t)) - \kappa_i(P_i(t))| \bigg\} \bigg) \bigg)
% \end{align}
% for all $t \geq 0$, where we used the triangle inequality for class $\mathcal{K}$ functions, $\breve{\rho}(x+y) \leq \breve{\rho}(2x) + \breve{\rho}(2y)$. Finally, with \eqref{eq:lemma6_Gamma_bar}, we arrive at
% \begin{align}
%     \Gamma(t) \leq& \breve{\rho} \bigg( 2\beta_1 \bigg( \rho \bigg( \Gamma(0) \bigg), t \bigg) \bigg) + \breve{\rho} \bigg( 2\gamma_1 \bigg( \sup_{0 \leq \tau \leq t} \bigg\{ \sum_{i=1}^m |\kappa_i(\hat{P}_i(t)) - \kappa_i(P_i(t))| \bigg\} \bigg) \bigg) \nonumber \\
%         \leq& \beta(\Gamma(0), t) + \alpha_3 \bigg( \sup_{0 \leq \tau \leq t} \bigg\{ \sum_{i=1}^m |\kappa_i(\hat{P}_i(t)) - \kappa_i(P_i(t))| \bigg\} \bigg) \nonumber \\
%         \leq& \beta(\Gamma(0), t) + \alpha_3 \bigg( \sum_{i=1}^m \sup_{0 \leq \tau \leq t} \bigg| \kappa_i(\hat{P}_i(t)) - \kappa_i(P_i(t)) \bigg| \bigg)
% \end{align}
% with $\beta(s,t) = \breve{\rho}(2\beta_1(\rho(s),t))$, and $\alpha_3(s) = \breve{\rho}(2\gamma_1(s))$. Now, noting that $\kappa_i$ is continuous (Assumption \ref{ass:Assumption5}), there exists a class $\mathcal{K}_\infty$ function $\alpha_4$ such that $\sup_{0 \leq \tau \leq t} |\kappa_i(\hat{P}_i(t)) - \kappa_i(P_i(t))| \leq \alpha_4(\epsilon)$, for all $i \in [m]$, where $\epsilon$ is as in Theorem \ref{thm:predictor_operator_theorem}. Thus, we have
% \begin{align}
%     \Gamma(t) \leq& \beta(\Gamma(0), t) + \alpha_3 \bigg( \sum_{i=1}^m \alpha_4(\epsilon) \bigg) \nonumber \\
%         \leq& \beta(\Gamma(0), t) + \alpha_1( m\epsilon),
% \end{align}
% for all $t \geq 0$, which completes the proof.
% \end{proof}
% Similarly to \cite{Luke_Peijia:2025}, the constraint on the initial condition is a direct consequence of the existence of neural operator approximation on compact intervals. Fortunately, one can make the radius of initial conditions $\Omega$ arbitrary large by increasing $\overline{\mathcal{B}}$. From practical standpoint, the increase of $\overline{\mathcal{B}}$ does not come without a cost. Since the range on which the neural operator approximation has to hold increased, one will need a larger training set and more neurons in the architecture to achieve the same value of the approximation error $\epsilon$.

% The interesting consequence of this result is that $\epsilon^*$ grows with $\overline{\mathcal{B}}$. Even though it may be counterintuitive, one can recognize that the space of all possible states increases with $\overline{\mathcal{B}}$, and, therefore, enables us to have a coarser approximation of the neural operator predictor. For example, if the neural operator approximation held globally, we would not need to worry about $\epsilon^*$, since no matter how large the approximation error is, we would always end up in a state compatible with the neural operator.

% \section{Numerical Results}
% \subsection{Mathematical background}
% We conduct a numerical experiment on an unicycle whose dynamics can be described by the following system
% \begin{align}
%     \dot{X}_1 =& U_2(t-D_2)\cos(X_3(t)) \label{eq:unicycle_x1_dot} \\
%     \dot{X}_2 =& U_2(t-D_2)\sin(X_3(t)) \label{eq:unicycle_x2_dot} \\
%     \dot{X}_3 =& U_1(t-D_1), \label{eq:unicycle_x3_dot}
% \end{align}
% where $(X_1, X_2) \in \mathbb{R}^2$ is the position of the unicycle, $X_3 \in \mathbb{R}$ is spatial orientation, $U_1 \in \mathbb{R}$ is the turning rate, and $U_2 \in \mathbb{R}$ is linear speed. The unicycle is a classic example of a nonholonomic system due to its constrained motion, which prevents instantaneous sideways movement. This constraint makes traditional static feedback control insufficient, necessitating a time-varying controller to achieve stabilization. Therefore, we consider the following predictor-feedback control law designed in \cite[Section 8.3]{Nonlinear_Control_Under_Nonconstant_Delay}
% \begin{align}
%     U_1(t) =& -5M_1(t)^2 \cos(3(t+D_1)) - M_1(t)X_1(t)\bigg(1 + 25\cos^2(3(t+D_1))\bigg) - P_{1,X_3}(t) \label{eq:unicycle_u1} \\
%     U_2(t) =& -M_2(t) + 5X_2(t)\bigg(\sin(3(t+D_2)) - \cos(3(t+D_2))\bigg) + X_2(t)U_1(t), \label{eq:unicycle_u2}
% \end{align}
% where
% \begin{align}
%     M_i(t) = P_{i,X_1}(t)\cos(P_{i,X_3}(t)) + P_{i,X_2}(t)\sin(P_{i,X_3}(t)) \label{eq:unicycle_P}, \quad i=1,2, \\
%     Q_i(t) = P_{i,X_1}(t)\sin(P_{i,X_3}(t)) - P_{i,X_2}(t)\cos(P_{i,X_3}(t)) \label{eq:unicycle_Q}, \quad i=1,2.
% \end{align}
% This controller achieves global asymptotic stabilization when $D_1 = D_2 = 0$. Therefore, in order to compensate for non-zero delays, we use predictors from \cite{Bekiaris-Liberis_Krstic:2016} given by
% \begin{align}
%     P_{1,X_1}(t) =& X_1(t) + \int_{t-D_1}^t U_2(\theta - D_{2,1}) \cos \bigg( P_{1,X_3}(\theta) \bigg) d\theta \label{eq:unicycle_P1X1} \\
%     P_{1,X_2}(t) =& X_2(t) + \int_{t-D_1}^t U_2(\theta - D_{2,1}) \sin \bigg( P_{1,X_3}(\theta) \bigg) d\theta \label{eq:unicycle_P1X2} \\
%     P_{1,X_3}(t) =& X_3(t) + \int_{t-D_1}^t U_1(\theta) d\theta \label{eq:unicycle_P1X3}
% \end{align}
% and
% \begin{align}
%     P_{2,X_1}(t) =& P_{1,X_1}(t) + \int_{t-D_{2,1}}^t U_2(\theta) \cos \bigg( P_{2,X_3}(\theta) \bigg) d\theta \label{eq:unicycle_P2X1} \\
%     P_{2,X_2}(t) =& P_{1,X_2}(t) + \int_{t-D_{2,1}}^t U_2(\theta) \sin \bigg( P_{2,X_3}(\theta) \bigg) d\theta \label{eq:unicycle_P2X2} \\
%     P_{2,X_3}(t) =& P_{1,X_3}(t) + \int_{t-D_{2,1}}^t \kappa_1(\theta + D_2, P_{2,X_1}(\theta), P_{2,X_2}(\theta), P_{2,X_3}(\theta)) d\theta, \label{eq:unicycle_P2X3}
% \end{align}
% where for all $t - D_{2,1} \leq \theta \leq t$
% \begin{align}
%     \kappa_1(\theta + D_2, P_{2,X_1}(\theta), P_{2,X_2}(\theta), P_{2,X_3}(\theta)) =& -M_2(\theta)^2 \cos (\theta + D_2) - M_2(\theta) X_2(\theta) \nonumber \\
%         &\qquad \times (1 + \cos^2 (\theta + D_2)) - P_{2,X_3}(\theta). \label{eq:unicycle_kappa}
% \end{align}

% \subsection{Dataset construction}
% In order to approximate predictor feedback, we developed several approaches such as successive approximation \cite{Predictor_Feedback_For_Delay_Systems} and classical neural operators (FNO \cite{FNO:2021} and DeepONet \cite{deeponet:2021}). It is important to emphasize that we chose only those ML models that satisfy the universal approximation theorem \ref{thm:predictor_operator_theorem}; otherwise, we would not have required theoretical guarantees.

% When it comes to the dataset for predictor operator training, we generate all the data using the successive approximation method. First, we simulated $100$ trajectories of length $T = 10\,\text{s}$ using a time step-size $\Delta t = 0.001$, where each trajectory starts in the same initial condition $(x, y, \theta) = [1, 1, 1]$ with uniform noise between $-0.2$ and $0.2$ to ensure the uniqueness of all trajectories. Furthermore, we assume that the initial control input history is zero, that is, $U_i(t) = 0$, $t \in [-D_i,0)$ for all $i \in [m]$.

% To train the neural operators, we performed a hyperparameter search for the optimization parameters (epochs, learning rate, weight decay, etc.) while using AdamW \cite{AdamW} as our optimizer with an exponential learning rate scheduler, as is standard practice for neural operators. We train using the standard $L_2$ loss and refer the reader to the Jupyter notebook example on GitHub \cite{ProjectGitHub} for full parameter details.

% \subsection{Results}
% We choose $D_1 = 0.25$ and $D_2 = 0.6$ for our demonstrative example. In Table \ref{tab:performance}, we present the performance of the neural approximated predictors from Definition \ref{def:delay_dependent_predictor_operators} throughout both training and validation datasets corresponding to time step $\Delta t=0.01$. All approaches achieve strong performance with errors on the magnitude of $10^{-5}$. Furthermore, in order to accurately evaluate the neural approximate predictor operator in a dynamical environment, we simulated 25 new trajectories for each approximation approach and calculated their average trajectory $L_2$ norm with respect to the stationary equilibrium point, $(x,y,\theta)=(0,0,0)$, and presented in column four of Table \ref{tab:performance}. Additionally, all initial conditions for those testing trajectories are sampled from the same distribution used for generating training datasets, which means that we neglected the $\epsilon$ term in \eqref{eq:Omega} due to high model accuracy.

% The right part of the table highlights the computational advantages of FNO and DeepONet over successive approximations, particularly for smaller time steps. As the time step decreases, the computational gain of these methods becomes pronounced. Successive approximations exhibit a linear scaling in computation time with respect to $\Delta t$, as seen in the increasing values from 0.379 ms ($\Delta t=0.01$) to 7.517 ms ($\Delta t=0.0005$). In contrast, FNO's computation time remains nearly constant across all time steps (e.g., ~3.8 ms), as its model is independent of the time step, requiring the same computational effort regardless of temporal resolution. On the other hand, DeepONet shows non-linear scaling, with computation times fluctuating between 1.024 ms and 3.842 ms. This behavior enables FNO and DeepONet to achieve computational advantages for finer time steps, where traditional methods like successive approximations become increasingly expensive. The consistency of FNO's performance underscores its efficiency in handling high-resolution temporal simulations without additional computational cost as illustrated in Figure \ref{fig:single_trajectory}.

% In summary, we briefly discuss the hurdles of designing a neural approximated predictor operator with distinct constant delays as well as the advantages of this approach. First, we observe that as the gap between delays increases, the amount of data necessary to achieve the same approximation accuracy increases. This is to be expected since the neural network has to learn multiple dynamics that correspond to the period between two consecutive control hits. On the other hand, the advantage of neural operator predictors over the traditional successive approximations approach becomes more pronounced in more challenging settings. One such example is presented in \cite{Luke_Peijia:2025}, where the cost of executing manipulator dynamics is significantly greater than the cost of a single forward pass in a neural network. This presents the most important trade-off in applying neural networks to control problems. One will have to spend more time on development, and it will be paid off through computational efficiency. In the end, such compromises are naturally unavoidable when implementing real-time controllers in complex systems.

% \begin{table}[htbp]
% \centering
% \caption{
%   The performance of predictor approximation approaches across training and validation datasets, as well as the average trajectory $L_2$ error across the testing dataset, with constant and distinct delays for unicycle stabilization with time step $\Delta t=0.01$ (left) and computation times for various step sizes (right).
% }
% \label{tab:performance}
% \setlength{\tabcolsep}{4pt}
% \resizebox{\textwidth}{!}{
% \begin{tabular}{@{}*{9}{c}@{}}
%     \toprule
%     \multirow{2}{*}{\makecell[c]{Approximation\\Method}} & 
%     \multirow{2}{*}{\makecell[c]{Parameters}} & 
%     \multirow{2}{*}{\makecell[c]{Training $L_2$\\error $(\times10^{-6})\downarrow$}} & 
%     \multirow{2}{*}{\makecell[c]{Validation $L_2$\\error $(\times10^{-6})\downarrow$}} & 
%     \multirow{2}{*}{\makecell[c]{Avg. trajectory\\$L_2$ error $\downarrow$}} & 
%     \multicolumn{4}{c}{\makecell[c]{Computation time (ms)}} \\
%     \cmidrule(lr){6-9}
%     & & & & & \makecell[c]{dx = 0.01} & \makecell[c]{dx = 0.005} & \makecell[c]{dx = 0.001} & \makecell[c]{dx = 0.0005} \\
%     \midrule
%     \makecell[c]{Successive\\approximations} & \makecell[c]{---} & \makecell[c]{---} & \makecell[c]{---} & 4.633 & 8.762 & 17.735 & 97.676 & 218.740 \\
%     FNO & 42211 & 0.33 & 0.25 & 5.121 & 30.850 & 29.664 & 34.653 & 31.093 \\
%     DeepONet & 1124824 & 3.00 & 2.46 & 4.420 & 7.245 & 9.704 & 10.816 & 23.126 \\
%     \bottomrule
% \end{tabular}
% }
% \end{table}

% \begin{figure}[htbp]
%     \centering
%     \includegraphics{sections/section3/plots/single_trajectory_FNO.png}
%     \caption{Example of a single trajectory of the unicycle system with the numerical and FNO approximated predictors}
%     \label{fig:single_trajectory}
% \end{figure}

\section{Conclusion and Future Work}
In conclusion, we extended the neural operator predictor framework to multi-input nonlinear systems with distinct delays, proving existence and semiglobal practical stability (Theorems \ref{thm:neural-op-approximate-predictor-theorem} and \ref{thm:main_result}). Our results, validated by unicycle simulations and consistent with single-delay cases \cite{Luke_Peijia:2025}, highlight the balance between efficiency and accuracy and applies to any approximate predictor satisfying uniform error bounds. This opens the door for more complex time-varying approximate predictor methodologies such as those with noise and output feedback designs. 
% \setlength{\bibitemsep}{0.8em}
\bibliography{references}
\bibliographystyle{abbrv}
% Appendix
% \appendix
% \renewcommand{\thesubsection}{\Alph{section}.\arabic{subsection}}
% \renewcommand{\thesection}{Appendix \Alph{section}:}


\end{document}
